\subsection{结论与证据类型}

本节在代数闭特征零域上证明下述定理；任意特征零域的结论由
Nullstellensatz 下降得到。上界是 Glynn 恒等式；下界由特征零代数
几何、Petersen 乘积影子、带符号小图割、有限整数表以及一个严格模三
端点证书组成。定理~\ref{n5:thm:one-intersection-flag} 的
\(886{,}464\) 个坐标旗标穷尽是逻辑前提；其余有限表的集合、递推、
情形数及逐项核查规则均在正文中证明，参考程序只核对抄录。此前唯一需要
\(\binom{15}{10}=3003\) 项精确有理数分类的
orbit--1 临界十平面，现由五顶点图闭式、十类非零 Möbius 项族和八行
\(K_4\) 表完成正文分类。因此五阶 lower--16 是特征零代数几何与有限
组合分类相结合的计算机辅助证明；除上述独立端点证书外，它不依赖其他
有限域秩、随机试验、SAT/DRAT、Gröbner 基或未公开数据。

\begin{theorem}\label{n5:thm:main}
在任意特征零域 \(k\) 上，
\[
  \ChowRank_k(\perm_5)=16.
\]
\end{theorem}

\subsection{导数空间、Koszul 映射与单项上界}

写 \(V=A\otimes B\)，其中
\(A=\langle a_0,\ldots,a_4\rangle\)、
\(B=\langle b_0,\ldots,b_4\rangle\)，并以
\(x_{ij}=a_i\otimes b_j\) 为变量。令
\[
 E=\im C_{3,2}(\perm_5),\qquad
 D=\im C_{2,3}(\perm_5).
\]
则
\begin{align}
 E&=\Span\{x_{ia}x_{jb}+x_{ib}x_{ja}:i<j,\ a<b\},\label{n5:eq:E}\\
 D&=\Span\{\perm(X_{I,J}):|I|=|J|=3\}.
\end{align}
这些生成元支撑互异，故
\[
 \dim E=\binom52^2=100,\qquad
 \dim D=\binom53^2=100.
\]

对二次空间 \(F\subset\Sym^2V\)，置
\[
 F^{(1)}=\{g\in\Sym^3V:\partial_vg\in F\text{ 对所有 }v\in V^*\}.
\]
Koszul 微分
\[
 \delta_F:V\otimes F\longrightarrow\Lambda^2V\otimes V,
 \qquad w\otimes uv\longmapsto
 u\wedge w\otimes v+v\wedge w\otimes u
\]
满足自然同构 \(\ker\delta_F\simeq F^{(1)}\)。固定三行三列后，六个
匹配单项式的换位图连通，故 \(E^{(1)}=D\)。因此
\begin{equation}\label{n5:eq:Kperm}
 \rank K(\perm_5)=25\cdot100-100=2400.
\end{equation}

若 \(T=\ell_1\cdots\ell_5\) 是一个 Chow 项，则其二次导数空间
\(F_T\) 由十个两因子积张成，其三次延拓至少包含十个三因子积。由
退化的上半连续性，
\begin{equation}\label{n5:eq:single240}
 \dim F_T\le10,\qquad \rank K(T)\le25\cdot10-10=240.
\end{equation}

\subsection{低秩二次型与一维交引理}

\begin{lemma}[秩不超过五的分类]\label{n5:lem:lowrank}
若 \(0\ne q\in E\) 且其对称矩阵秩不超过五，则
\[
 M(q)=U\otimes W,\qquad \rank U=\rank W=2.
\]
特别地，\(q\) 的秩为四。
\end{lemma}

\begin{proof}
按行把 \(M(q)\) 写成 \(5\times5\) 块矩阵 \((D_{ij})\)。每个
\(D_{ij}\) 对称且零对角，\(D_{ii}=0\)。取非零块 \(D\)。主块
\(\left[\begin{smallmatrix}0&D\\D&0\end{smallmatrix}\right]\) 的秩为
\(2\rank D\)，而非零对称零对角矩阵不可能秩一，故
\(\rank D=2\)。若任何连接块在该主块余核中有非零像，对称加边会使
总秩至少增加二，达到六。因此所有连接块的像均在 \(\im D\) 中。

秩二对称零对角矩阵可写成 \(ab^{\mathsf T}+ba^{\mathsf T}\)，其中
\(a,b\) 支撑不交。在固定像 \(\langle a,b\rangle\) 中，对称零对角
条件迫使任何矩阵均为 \(D\) 的标量倍。对初始两块作合同 Schur 消元；
剩余补矩阵秩至多一，且每个块仍对称零对角。非零对角块秩至少二，
非零非对角块连同其转置给出秩至少四，故 Schur 补为零。于是所有块
都是 \(D\) 的倍数，即 \(M(q)=U\otimes D\)。秩乘法和零对角条件给出
\(\rank U=\rank D=2\)。
\end{proof}

\begin{corollary}\label{n5:cor:intersection}
若 \(\dim L\le5\)，则
\[
 \dim(E\cap\Sym^2L)\le1.
\]
因而对任意五次 Chow 项 \(T\)，
\(\dim(E\cap F_T)\le1\)。
\end{corollary}

\begin{proof}
若交空间含一条射影线，则引理~\ref{n5:lem:lowrank} 把这条线送入
秩二张量的 Segre 簇。完全包含在 Segre 簇中的射影线必须固定一个
因子。不妨固定左侧秩二像。两个不成比例的右因子具有不同的二维像，
其像之和至少三维，故本质变量空间至少为 \(2\cdot3=6\)，与
\(\dim L\le5\) 矛盾。
\end{proof}

\subsection{相对延拓与小维全局界}

令 \(Q=\Sym^2V/E\)。对 \(W\subset Q\)，记
\[
 p(W)=\dim(E+\widetilde W)^{(1)}-\dim D,
\]
其中 \(\widetilde W\) 是任意提升；该数与提升无关。行列环面把任意
Grassmann 点平坦退化到坐标点，核维在特殊纤维只能上升。因此
\(p(W)\) 的全局上界可在 225 个环面权组成的坐标集合上证明。这些权
分成 25 个平方、50 个同行边、50 个同列边和 100 个矩形 crossing。

这里同时固定局部系数图的规范。对一个三次行列环面权 \(\omega\)，
以该权的三次单项式为顶点。对每个变量 \(x_{ia}\) 和二次商权 \(q\)，
条件 \(\partial_{ia}f\in E+\widetilde W\) 给出一条系数行：若该行有
两个非零项，就画一条带符号的 \(q\)-标号边；若只有一个非零项，就把
该顶点以 \(q\) 锚定到零。选择 \(q\in W\) 正是删除全部 \(q\)-标号
关系。下述四边形、六边图及 \(K_{3,3}\) 块的符号在换顶点号后相容，
故其局部核维数等于不含锚点的连通分量数。于是加入一个新权 \(q\) 的
边际，就是删除 \(q\)-关系后新产生的未锚连通分量数。这个定义完全在
整数系数 \(\{0,\pm1\}\) 上进行。

\subsection{一维交旗标的普适上界}

记 \(\pi:\Sym^2V\to Q\) 为商映射。下述结果补足从
\(\dim(E\cap F_T)=1\) 到相对延拓界之间所需的非平凡桥梁；其中有限
端点的穷尽是定理的逻辑前提，而不是冗余诊断。

\begin{theorem}[一维交旗标定理（计算机辅助）]
\label{n5:thm:one-intersection-flag}
设 \(T=\ell_1\cdots\ell_5\) 为五次 Chow 项，且
\[
 \dim F_T=10,\qquad \dim(E\cap F_T)=1.
\]
令 \(Z=\pi(F_T)\)，则 \(\dim Z=9\)。若
\[
 Z\subset W\subset Q,\qquad \dim W=10,
\]
则 \(p(W)\le26\)。
\end{theorem}

\begin{proof}
先扩张到代数闭包。令
\(L_T=\langle\ell_1,\ldots,\ell_5\rangle\)，并置
\(d=\dim L_T\)。因为
\(F_T\subset\Sym^2L_T\)、\(\dim F_T=10\) 且 \(d\le5\)，故
\(d\in\{4,5\}\)。令 \(\mathbb T\) 为行列对角环面。对固定的 \(d\)，
考虑下列射影乘积中的闭关联簇：
\[
 \mathcal I_d\subset
 \operatorname{Gr}(d,V)\times\operatorname{Gr}(10,\Sym^2V)
 \times\operatorname{Gr}(9,Q)\times\operatorname{Gr}(10,Q)
 \times\mathbf P(E),
\]
其点 \((L,F,Z',W',[q])\) 满足
\begin{equation}\label{n5:eq:oneflag-incidence}
 F\subset\Sym^2L,\qquad q\in F,\qquad
 F\subset\pi^{-1}(Z'),\qquad Z'\subset W'.
\end{equation}
这些条件都是相应万有子丛之间自然映射的零化条件，因而
\(\mathcal I_d\) 射影且 \(\mathbb T\)-稳定。

取 \(0\ne q\in E\cap F_T\)。原数据给出
\((L_T,F_T,Z,W,[q])\in\mathcal I_d\)。其 \(\mathbb T\)-轨道闭包是
完备的；Borel 不动点定理给出一个 \(\mathbb T\)-不动点
\[
 (L_0,F_0,Z_0,W_0,[q_0]).
\]
推论~\ref{n5:cor:intersection} 与
\(0\ne q_0\in E\cap F_0\subset E\cap\Sym^2L_0\) 给出
\begin{equation}\label{n5:eq:oneflag-rankpreserve}
 E\cap F_0=\langle q_0\rangle,\qquad
 \dim\pi(F_0)=9,\qquad Z_0=\pi(F_0).
\end{equation}
最后一个等号来自
\(\pi(F_0)\subset Z_0\) 以及两边同为九维。特别地，这里没有把
“商像维数在特殊纤维保持”误当成闭条件；它由闭关联、固定的
Grassmann 维数和一维交推论共同推出。

下面分类这个不动端点。\(V\) 的 \(25\) 个坐标线是互异权空间，故
\(L_0\) 是坐标 \(d\)-平面。\(E\) 的 \(100\) 个互异权线为
\[
 \left\langle x_{ia}x_{jb}+x_{ib}x_{ja}\right\rangle,
 \qquad i<j,\quad a<b,
\]
所以 \(q_0\) 是一个矩形 permanent。该二次型的本质变量空间恰为
\[
 \langle x_{ia},x_{ib},x_{ja},x_{jb}\rangle,
\]
并包含于 \(L_0\)。若 \(d=4\)，则 \(L_0\) 正是这个矩形，且由维数
得到 \(F_0=\Sym^2L_0\)；其商像是该矩形的九个互异商权。

若 \(d=5\)，则 \(L_0\) 是矩形再加一个坐标格。模行列置换与转置，
只有两类：第五格与矩形共享一行或一列的 attached 类，以及既不共享行
也不共享列的 external 类。每个矩形分别有 \(12\) 与 \(9\) 种第五格，
故在全部坐标五平面中两类计数为
\begin{equation}\label{n5:eq:oneflag-orbits}
 100\cdot12=1200,\qquad 100\cdot9=900.
\end{equation}
由推论~\ref{n5:cor:intersection}，
\(\pi(\Sym^2L_0)\) 恰有 \(14\) 维；其十四个环面权均为一维。
式~\eqref{n5:eq:oneflag-rankpreserve} 表明 \(Z_0\) 是其中任意可能的
九权子集。另一方面，\(Q\) 的 \(225\) 个互异一维权空间分为
\(25\) 个 square、\(50\) 个同行边、\(50\) 个同列边和 \(100\) 个
crossing。因此 \(W_0\) 是 \(Z_0\) 的九个权再加 \(225\) 个权中的任意
一个新权；这个第十权并未被假设落在上述十四权宇宙中。

只剩一个有限而精确的端点命题。对坐标 \(W_0\)，令
\[
 \Phi_{W_0}:\Sym^3V\longrightarrow V\otimes(Q/W_0)
\]
为极化偏导后投影到 \(Q/W_0\) 的映射，则
\(p(W_0)=\dim\ker\Phi_{W_0}-100\)。在三次 divided-power 基和上述
商权基下，\(\Phi_{W_0}\) 是元素属于 \(\{0,\pm1\}\) 的整数矩阵。
因此
\[
 \rank_{\mathbf F_3}\Phi_{W_0}
 \le\rank_{\Q}\Phi_{W_0},\qquad
 \dim_{\Q}\ker\Phi_{W_0}
 \le\dim_{\mathbf F_3}\ker\Phi_{W_0}.
\]
记 \(\overline\Phi_{W_0}\) 为这个整数矩阵逐项模三后的矩阵，并定义
\[
 \overline p_3(W_0)=
 \dim_{\mathbf F_3}\ker\overline\Phi_{W_0}-100.
\]
独立重建程序
\path{perm5_one_intersection_flag_standalone_exact.py} 不读取任何冻结结果，
而从极化定义重建 \(225\) 个二次商权、\(2925\) 个三次 divided-power
单项式及 \(1225\) 个环面权块。它先严格得到模三基底矩阵
\(\overline\Phi_0\) 的核维为 \(100\)，与特征零中的
\(\dim_{\Q}D=100\) 相同，再穷尽下表中的端点：
\begin{equation}\label{n5:eq:oneflag-finite-table}
\begin{array}{c|r|c}
\text{因子张成维数与轨道}&\text{检查的旗标数}&
 \max \overline p_3(W_0)\\ \hline
4&100(225-9)=21{,}600&22\\
5\text{，attached}&\binom{14}{9}(225-9)=432{,}432&26\\
5\text{，external}&\binom{14}{9}(225-9)=432{,}432&22
\end{array}
\end{equation}
每个局部矩阵均以模三精确消元计算；附录
\ref{app:one-intersection-flag-spec} 给出对象、伪代码、完整直方图、失败
语义和输出哈希。故 \(p_{\Q}(W_0)\le26\)。

最后，\(p(tW)=p(W)\) 对一切 \(t\in\mathbb T\) 成立，因为 \(E\) 与
极化映射均为 \(\mathbb T\)-等变。\(\ker\Phi_W\) 又是 Grassmann 簇上
一个向量丛映射的核，故其维数上半连续。沿上述轨道特殊化得到
\[
 p(W)=p(tW)\le p(W_0)\le26.
\]
矩阵秩在扩域下不变，故结论从代数闭包下降到任意特征零域。
\end{proof}

正文所需的精确界为
\begin{equation}\label{n5:eq:p-small}
\begin{aligned}
 (p_0,\ldots,p_8)&=(0,1,4,5,8,11,20,24,28),\\
 p_9&\le35,\qquad p_{10}\le50,\qquad
 p_{11}\le55,\qquad p_{12}\le61.
\end{aligned}
\end{equation}
其中 \(p_d=\max_{\dim W=d}p(W)\)。先用纯压缩论证处理无 crossing
集合，再用闭式割密度加入 crossing。

\begin{lemma}[无 crossing 压缩与闭式包络]
\label{n5:lem:nocrossing}
若 \(F\) 是由 square、同行边和同列边组成的坐标权集合，且
\(|F|=d\)，则
\[
 p(F)\le35,50,55,60\qquad(d=9,10,11,12).
\]
此外，\(d=10\) 且含 square 时 \(p(F)\le43\)；\(d=11\) 时，除
“一层 \(K_5\) 加一个 square”外均有 \(p(F)\le50\)。
当 \(d=9\) 时，等号恰为一个固定 row 或 column 中的
\(K_5-e\)，故共有一百个等号集合。
\end{lemma}

\begin{proof}
把 square 集记为 \(A\subset[5]\times[5]\)。对每行 \(i\)，令
\(G_i\) 是列顶点上的同行边图；对每列 \(a\)，令 \(H_a\) 是行顶点
上的同列边图。由上面的局部系数图定义，逐个无锚分量计数得到
\begin{equation}\label{n5:eq:nocrossing-formula}
 p(A,G,H)=5\sum_i t(G_i)+5\sum_a t(H_a)
 +\sum_{(i,a)\in A}(d_{G_i}(a)+1)(d_{H_a}(i)+1).
\end{equation}
这里的分量可直接写出。每个 square \(S_{ia}\) 给出 \(x_{ia}^3\)；
它与一条端点为 \(a\) 的 row edge 给出 \(x_{ia}^2x_{ib}\)，column
情形对称。\(G_i\) 的每个三角形 \(\{a,b,c\}\) 给出行内单项式
\(x_{ia}x_{ib}x_{ic}\)，并对四个 \(j\ne i\) 分别给出
\[
x_{ia}x_{ib}x_{jc}+x_{ia}x_{ic}x_{jb}
+x_{ib}x_{ic}x_{ja},
\]
合计五维；column 情形对称。最后，若
\(S_{ia},R_{i;ab},C_{ij;a}\) 同时选中，则
\[
2x_{ia}x_{ib}x_{ja}+x_{ia}^2x_{jb}
\]
给出一个 corner。逐变量求导可见这些三次型都在相对延拓中；反过来，
按三次单项式的行、列重数分类，不含 crossing 时只能出现上述四类，
且不同环面权块线性无关。这就证明
式~\eqref{n5:eq:nocrossing-formula}。最后一项正是 square、
square--row-edge、square--column-edge 与 corner 四类模式计数之和。

先证可同时压缩所有行、列标号。固定 \(u<v\) 并压缩行标号；于是
square 的两个行切片作 singleton 压缩，
\(G_u,G_v\) 变成 \(G_u\cup G_v,G_u\cap G_v\)，每个 \(H_a\)
在顶点 \(u,v\) 上作标准图压缩。所需的二点重排恒等式是
\begin{equation}\label{n5:eq:binary-rearrangement}
 (x\vee y)(a\vee b)+(x\wedge y)(a\wedge b)\ge xa+yb
 \qquad(x,y,a,b\in\{0,1\}).
\end{equation}
若 \(X\) 是顶点集、\(K\) 是图，置
\(I(X,K)=\sum_{v\in X}d_K(v)\)。对每个
\(k\notin\{u,v\}\)，压缩前后边 \(uk,vk\) 对 \(I\) 的贡献之差
正是式~\eqref{n5:eq:binary-rearrangement} 的左、右差；边 \(uv\)
的贡献不变。因此同时压缩 \(X,K\) 不降低 \(I(X,K)\)。

现在逐项检查式~\eqref{n5:eq:nocrossing-formula}。square 数不变；
square--row-edge 关联直接用式~\eqref{n5:eq:binary-rearrangement}；
square--column-edge 关联是 \(I(A^a,H_a)\)。对 corner，固定
\(a\ne b\) 并令
\[
 X_{a,b}=\{i:(i,a)\in A,\ ab\in G_i\}.
\]
压缩后的中心集包含 \(C_{uv}X_{a,b}\)，故
\(I(X_{a,b},H_a)\) 不减。最后，三角形指标对边集超模，因而
\[
 t(K\cup L)+t(K\cap L)\ge t(K)+t(L);
\]
标准图压缩也不减少三角形：把含 \(v\) 不含 \(u\) 的三角形压到
\(u\)，若目标原已存在，则原三角形的两条 \(v\)-边都不会被删除。
所以行压缩不降低 \(p\)，列压缩完全对称。

先压行再压列；图压缩对包含关系单调，shifted 图的并、交仍 shifted，
故最终可设
\[
G_0\supseteq\cdots\supseteq G_4,\qquad
H_0\supseteq\cdots\supseteq H_4,
\]
每个图均 shifted，且 \(A\) 是 Ferrers ideal。

五个简单图共有 \(e\le12\) 条边时，Kruskal--Katona 定理给三角形
包络
\begin{equation}\label{n5:eq:triangle-envelope}
 \tau(e)=(0,0,0,1,1,2,4,4,5,7,10,10,10)_e.
\end{equation}
确切地，若 \(e=\binom{k}{2}+\ell\)、\(0\le\ell<k\)，单图的三角形
数至多 \(\binom{k}{3}+\binom{\ell}{2}\)；对 \(e\le10\) 合并边预算
只会增大该包络。对 \(e=11,12\)，一个 \(K_5\) 加一、二条别层边给
十个三角形；没有 \(K_5\) 时至多分别为 \(7+0,7+1\)。此外
\(\tau(6)=4,\tau(9)=7,\tau(10)=10\) 的等号形状分别是一层
\(K_4,K_5-e,K_5\)。

置 \(s=|A|\)、\(r=\sum_i|E(G_i)|\)、
\(c=\sum_a|E(H_a)|\)。令 \(R,C,L\) 分别是所有 selected square
上的 row 度和、column 度和、度数乘积和。每条纯边至多碰到两个
square，每个度至多四，故
\[
R\le R_0:=\min(4s,2r),\qquad
C\le C_0:=\min(4s,2c).
\]
而 \(L\) 计数一对 row/column edge 的共同 selected 中心；每对边至多
有一个中心，故 \(L\le rc\)，同时 \(L\le4R,4C\)。于是 square 项
\(Q\) 满足
\begin{equation}\label{n5:eq:universal-square-bound}
 Q\le U_s(r,c):=
 s+R_0+C_0+\min\{rc,4R_0,4C_0\}.
\end{equation}
结合式~\eqref{n5:eq:triangle-envelope}，
\begin{equation}\label{n5:eq:nocrossing-master}
 p\le5\bigl(\tau(r)+\tau(c)\bigr)+U_s(r,c),
 \qquad s+r+c=d.
\end{equation}

把 \(d=9,10,11,12\) 代入；按 \(r,c\) 转置只计一次，式
\eqref{n5:eq:nocrossing-master} 超过目标的整数情形恰为
\begin{center}
\begin{tabular}{@{}ccl@{}}
\toprule
\(d\)&目标&例外 \((s;r,c)\)\\ \midrule
9&35&\((1;2,6),(1;3,5),(2;1,6),(2;3,4)\)\\
10&50&无\\
11&55&\((2;3,6)\)\\
12&60&\((1;1,10),(2;3,7),(2;4,6),(2;5,5),(3;3,6)\)\\
\bottomrule
\end{tabular}
\end{center}
这只是把 \(r+c=d-s\) 代入显式式~\eqref{n5:eq:nocrossing-master}；
下面逐式闭合全部候选。

若 \(s=1\)，则
\begin{equation}\label{n5:eq:one-square-bound}
 Q\le(\min(4,r)+1)(\min(4,c)+1).
\end{equation}
三个例外分别给 \(20+15=35,15+20=35,50+10=60\)。九维的前两个
虚假等号还可严格排除：在 \((r,c)=(2,6)\) 中，四个三角形迫使六边
一侧为 \(K_4\)，两侧标记度至多二、三，故 \(p\le20+12=32\)；
三角形少于四时 \(p\le15+15=30\)。在 \((3,5)\) 中，三个三角形
迫使两侧极值图为 \(K_3,K_4-e\)，标记度至多二、三，故
\(p\le15+12=27\)；较少三角形时至多三十。最后 \((4,4)\) 的两个
三角形使两个标记度都至多三，给 \(p\le10+16=26\)；否则至多三十。
所以九维 \(s=1\) 时实际 \(p\le32\)。

若 \(s=2\)，转置后 Ferrers ideal 是 row domino
\(\{(0,0),(0,1)\}\)。写两格的 row 度和、column 度和为 \(G,H\)。
只有 row 图中的边 \(01\) 能被 \(G\) 重复计数，而 \(H\) 来自两个
column 图层，所以
\[
G\le\min\{8,r+\mathbf1_{\{r>0\}}\},\qquad
H\le\min\{8,c\},
\]
并且
\begin{equation}\label{n5:eq:domino-bound}
 Q\le2+G+H+\min\{rc,4G,4H\}.
\end{equation}
式~\eqref{n5:eq:domino-bound} 关闭表中所有 \(s=2\) 情形及其两个方向，
唯一尚超过目标的粗值是九维 row-domino 的 \((r,c)=(1,6)\)。这时
唯一 row edge 是 \(01\)。若六边一侧有四个三角形，它是一层
\(K_4\)，度向量为 \((g_0,g_1)=(1,1)\)、
\((h_0,h_1)=(3,0)\)，故 \(Q=10,p=30\)；否则三角形至多三且
式~\eqref{n5:eq:domino-bound} 给 \(Q\le16,p\le31\)。

若 \(s=3\)，转置后 Ferrers ideal 只有 row triple 与 L 形。前者
满足
\[
R\le\min\{12,r+\min(r,3)\},\qquad C\le\min\{12,c\};
\]
后者满足
\[
R\le\min\{12,r+\mathbf1_{\{r>0\}}\},\qquad
C\le\min\{12,c+\mathbf1_{\{c>0\}}\}.
\]
两者仍有 \(L\le rc\)。十二维唯一例外 \((r,c)=(3,6)\) 及其转置，
在两形中分别给 \(Q\le33,32\)；由
\(\tau(3)+\tau(6)=5\)，得到 \(p\le58,57\)。九维 \(r+c=6\)
同样逐式给 \(p\le32\)。

九维 \(s\ge4\) 时，式~\eqref{n5:eq:nocrossing-master} 已严格小于
三十五；所以九维等号只能有 \(s=0\)。此时 \(p=5T\)，而
\(T=\tau(9)=7\) 的等号形状恰是一层 \(K_5-e\)。共有五个 row、
五个 column，且每层有十种缺边，故等号集合数为一百。

最后记录后面要用的两个余量。\(d=10,s=1\) 时，式
~\eqref{n5:eq:one-square-bound} 只有 \((r,c)=(3,6)\) 及其转置给出
粗值四十五；五个三角形迫使两侧为 \(K_3,K_4\)，标记度至多二、三，
故实际至多 \(25+12=37\)，其余拆分至多四十。式
~\eqref{n5:eq:domino-bound} 对 \(s=2\) 至多给四十三，三格两形至多
三十九，\(s\ge4\) 的式~\eqref{n5:eq:nocrossing-master} 至多三十六，
故含 square 时 \(p\le43\)。

\(d=11\) 时，\(s=1\) 的式~\eqref{n5:eq:one-square-bound} 只有
\((r,c)=(0,10),(10,0)\) 达到五十五，恰为一层 \(K_5\) 加一个
square；其余至多五十。\(s=2\) 的 domino 粗界只有
\((r,c)=(3,6)\) 可能给五十三；五个三角形时两侧仍为 \(K_3,K_4\)，
直接代入两个 domino 度向量给 \(Q\le16,p\le41\)，少一个三角形时
由式~\eqref{n5:eq:domino-bound} 得 \(p\le48\)。三格两形至多四十六，
\(s\ge4\) 至多四十四。这证明所述余量。
\end{proof}

\begin{lemma}[crossing 边际密度]\label{n5:lem:crossdensity}
若已有坐标方向数为 \(N\le11\)，加入一个 crossing 权 \(x\) 后，
\[
 \Delta_xp\le
 \begin{cases}
 N,&0\le N\le4,\\
 \lceil3N/4\rceil,&5\le N\le11.
 \end{cases}                                                   \tag{CM}
\]
\end{lemma}

\begin{proof}
固定 \(x=X_{01;01}\)。免费加入其四条矩形边界后，十二个重复块分为
六组，九个 matching 块组成 \(3\times3\) 格。每个重复组有两条
side crossing；组内两个六边块各自还需要一条属于自己的 external
edge 才能隔开 \(x\) 的两端，且两块的 external 标签互不相交。因此
激活一个重复块至少花一条 external edge，并要求该组 side-positive；
一个组至多承载两个重复块。

对一个 matching 格，六个顶点是相应三行三列的六个完美匹配；奇偶
二分后关系图是 \(K_{3,3}\)，边以交换所用的 \(2\times2\) 矩形标号。
除 \(x\) 外，四条 side 标签分属该格的横、纵两组，剩余四条
diagonal 标签只属于该格。把六个置换逐个列出即得：
\begin{enumerate}[label=(\roman*),leftmargin=2.4em]
\item 横、纵两组都 double-side 时，diagonal 代价可为零；
\item 两组都 side-positive 时，除前一种情形外至少需一条 diagonal；
\item 其余情形至少需两条 diagonal。
\end{enumerate}
不同格的 diagonal 标签互不相交。以上是六边图与 \(K_{3,3}\) 的完整
割输入；下面只做整数汇总。令
\[
\begin{aligned}
u&=h_2+h_1,& v&=v_2+v_1,&
\alpha&=h_2,& \gamma&=v_2,\\
P&=u+v,& A&=P+\alpha+\gamma,&
z&=\alpha\gamma,& o&=uv.
\end{aligned}
\]
这里 \(h_i,v_i\) 计数选择 \(i\) 条 side 的横、纵组，故
\(0\le\alpha\le u\le3\)、\(0\le\gamma\le v\le3\)。
若激活 \(r\) 个重复块、\(m\) 个 matching 块，则六边图与
\(K_{3,3}\) 的最小割分别给
\[
r\le2P,\qquad
q\ge Q:=r+A+D(m;z,o),                                      \tag{CM1}
\]
其中
\[
D(m;z,o)=
\begin{cases}
0,&m\le z,\\
m-z,&z<m\le o,\\
2m-o-z,&m>o.
\end{cases}                                                   \tag{CM2}
\]
三项依次对应 zero-cost、one-cost 和 two-cost 的独有 diagonal。

若 \(Q\le11\)，则
\begin{equation}\label{n5:eq:density}
4(r+m)\le3Q.
\end{equation}
为证此式，先注意 \(z>0\) 时
\[
z+3\le2(\alpha+\gamma)\le A.                                \tag{CM3}
\]
又若 \(A\le9\)，则
\[
o+3z\le P+3\alpha+3\gamma+2,                                \tag{CM4}
\]
唯一例外是 \((u,v,\alpha,\gamma)=(3,3,0,0)\)。事实上令
\(w=\alpha+\gamma\)；由 \(A\ge2w\) 得 \(w\le4\)。
\(w=0\) 时只可能在 \(u=v=3\) 失败；\(1\le w\le3\) 时使用
\(uv-u-v\le3\) 与
\(\alpha\gamma-\alpha-\gamma\le-1\)；\(w=4\) 时
\(u+v\le5\)，两项分别至多一、零。

在 \(m\le z\) 段，若 \(m>0\)，由预算与 (CM3)
\[
r+4m\le11-A+4z\le3A-1;
\]
\(m=0\) 则由 \(r\le2P\le2A\) 直接成立。在
\(z<m\le o\) 段，所求等价于
\(r+m+3z\le3A\)：\(z>0\) 仍用同一行；\(z=0,A\ge3\) 时左端至多八；
\(A\le2\) 时只能 \(u=v=1,A=2\)，故 \(r+m\le5\)。
最后 \(m>o\) 时，\(A\ge10\) 会由 \(D\ge2\) 违反预算；否则 (CM4)
与 \(m\ge o+1\) 给
\[
r+3o+3z\le2P+3o+3z\le3A+2o+2\le3A+2m.
\]
例外型有 \(o=9\)，不可能再取 \(m>o\)。这证明
\eqref{n5:eq:density}，从而非角块数满足
\[
F(q)\le\lfloor3q/4\rfloor\qquad(0\le q\le11).                \tag{CM5}
\]
当 \(q=4\) 还有 \(F(4)\le2\)：若三块被激活，按 (CM2) 三段，
\(m\le z\) 中若 \(m=0\)，则 \(r\ge3,P\ge2\) 给 \(Q\ge5\)；
若 \(m>0\)，则 \(A\ge4\)，预算迫使
\(u=v=\alpha=\gamma=1,r=0,m\le1\)。
\(z<m\le o\) 中，\(z>0\) 给 \(A+D\ge5\)；\(z=0\) 给
\(A\ge2,D=m\)，故 \(r+m\le2\)。
最后 \(m>o\) 时 \(D\ge2\)。若 \(r>0\)，预算唯一可能为
\(A=r=1,D=2\)，从而 \(o=z=0,m=1\)；若 \(r=0,m\ge3\)，则
\(o\ge2\)、\(o=0\)、\(o=1\) 三种情形分别给
\(A+D\ge5,D\ge6,A+D\ge6\)。均矛盾。

现在把 \(N\) 个已有方向分成 \(s\) 个角 square、\(b\) 条边界边和
\(q=N-s-b\) 个非角方向。边界中横、纵边数之积至多
\[
g(b)=\lfloor b^2/4\rfloor=(0,0,1,2,4)_b,
\]
故角块数 \(c\le\min\{4,s+g(b)\}\)。若 \(N\le4\)，
\(c\le s+b,F(q)\le q\)，即得 \(c+F(q)\le N\)。
若 \(5\le N\le11\)，置 \(a=s+b\)；直接按 \(b=0,\ldots,4\)
比较上式得到
\[
c\le\lceil3a/4\rceil,
\]
仅 \((s,b)=(4,0),(0,4)\) 例外：\(b=0\) 只需分 \(s=4\)；
\(b=1\) 用 \(4s\le3s+3\)（\(s=4\) 直接代入）；\(b=2\) 分
\(s\le2,s\ge3\)；\(b=3\) 分 \(s=0,1,s\ge2\)；\(b=4\) 仅
\(s=0\) 例外。非例外时与 (CM5) 相加即得 (CM)。
例外时 \(q=N-4\in[1,7]\)；若 \(q\ne4\)，
\(\lfloor3q/4\rfloor=\lceil3q/4\rceil-1\)，而 \(q=4\) 使用
\(F(4)\le2\)，结论仍同。
\end{proof}

\begin{corollary}[完整纯边线的 crossing 边际]\label{n5:cor:full-line}
若十个 no-crossing 方向是一层 pure row 或 column \(K_5\)，则加入任意
crossing 的边际为零；若再已有一个 square，则加入 crossing 的边际
至多一。
\end{corollary}

\begin{proof}
回看引理~\ref{n5:lem:crossdensity} 的完整局部割分解。pure line 不含
square 或 side crossing；在 \(x\) 的四条矩形边界中，它只占一个方向
的一条 boundary，因而不能形成横纵 boundary 对。故四个 corner 块
均不激活。每个六边 repeat 块要求一个 side-positive 组与自己的
external edge。matching 块若没有横纵 side-positive，则第三种割还
必须选择两条 diagonal crossing 标签。pure line 虽可提供 external
标签，却既不提供 side-positive 组，也不含 diagonal crossing 标签。
因此所有非角块为零。增加
一个 square 不改变非角块；它至多是 \(x\)
四个 corner cell 中的一个，故至多新激活一个 corner 块。
\end{proof}

现在完整推出式~\eqref{n5:eq:p-small}。把
引理~\ref{n5:lem:crossdensity} 的边际上界记为
\[
(\mu_0,\ldots,\mu_{11})=(0,1,2,3,4,4,5,6,6,7,8,9).
\]
九维集合若含 crossing，最后加入一个 crossing 给
\(p\le p_8+\mu_8=28+6=34\)；否则用
引理~\ref{n5:lem:nocrossing}。十维含 crossing 时
\(p\le p_9+\mu_9=42\)，无 crossing 时至多五十。

十一维无 crossing 时至多五十五。若恰有一个 crossing，去掉它后的
十维集合若含 square，则由引理~\ref{n5:lem:nocrossing} 的余量至多
\(43+8=51\)；若不含 square 且不是 \(K_5\)，十条纯边的三角形总数
至多七，故至多 \(35+8=43\)；若是 \(K_5\)，用
推论~\ref{n5:cor:full-line} 得五十。若至少有两个 crossing，把它们
全部移去并按 \(\mu_N\) 加回；对 crossing 数
\(2,\ldots,11\) 的最大粗值为五十二。因此 \(p_{11}\le55\)。

十二维无 crossing 时至多六十。恰有一个 crossing 时，十一维 base
若是一层 \(K_5\) 加一个 square，则推论~\ref{n5:cor:full-line} 给
\(55+1=56\)；其余 base 至多五十，故至多 \(50+9=59\)。
恰有两个 crossing 时，十维 base 若为 \(K_5\)，先加者边际为零、后加者
至多九，得五十九；base 含 square 时至多
\(43+8+9=60\)；不含 square 且不是 \(K_5\) 时至多
\(35+8+9=52\)。最后，crossing 数 \(3,\ldots,12\) 的逐层粗值为
\[
59,58,60,61,56,57,57,58,56,55,
\]
最大为六十一。这证明式~\eqref{n5:eq:p-small}。

旧 square 轨道表、按 square 数的无 crossing 最大值表和十二项
crossing 边际表均不再承担证明责任；附录中的整数程序只作公式反例
筛查，不使用任何域上的矩阵秩。

\subsection{不经 lower--15 的 fixed-six 入口}

式~\eqref{n5:eq:Kperm}--\eqref{n5:eq:single240} 先给出
\[
  \ChowRank(\perm_5)\ge\left\lceil\frac{2400}{240}\right\rceil=10.
\]
反设 Chow 秩至多十五，并取一个最小的 (r) 项分解；于是
\(10\le r\le15\)。固定其中六个原始项。对余下任一非零 Chow 项
\(T\)，恒等式
\[
 T=\tfrac12T+\tfrac12T
\]
把它拆成两个非零 Chow 项，并使项数增加一。反复拆分，便得到
\[
 \perm_5=T_1+\cdots+T_6+T_7+\cdots+T_{15},
\]
其中前六项仍是最小分解中的原始项，而余式
\(R=T_7+\cdots+T_{15}\) 恰为九个 Chow 项之和。因此下述 fixed-six
反证直接排除所有 \(r\le15\)，逻辑上不需要先建立 lower--15，也不
调用旧的 lower--15 SAT/DRAT 层。

令前六项之和为 \(C\)，并置
\[
 J=\im C_{2,3}(C),\quad H=\im C_{3,2}(C),\quad
 S=D\cap J,\quad T=E\cap H,
\]
以及
\[
 s=\dim S,\quad t=\dim T,\quad d=\dim(H/T),\quad h=t+d.
\]
记六个单项二次导数空间为 \(F_i\)，\(U=\sum_iF_i\)。注意
\(H\subset U\)，一般不能把同一催化矩阵的像 \(H\) 误写成六个像的
形式和 \(U\)。

三元 permanent 基以三行集、三列集标记。将三元集换成其补边后，
偏导影子成为 Petersen 图 \(P\) 的直积 \(P\times P\) 中的开邻域。
若 \(r=0,\ldots,23\)，十层动态递推给出最小影子
\begin{equation}\label{n5:eq:shadow-table}
\begin{array}{c|rrrrrrrrrrrr}
r&0&1&2&3&4&5&6&7&8&9&10&11\\ \hline
m_r&0&9&15&18&18&24&27&27&30&30&30&35\\[1mm]
r&12&13&14&15&16&17&18&19&20&21&22&23\\ \hline
m_r&35&36&36&36&36&42&45&45&48&48&48&52.
\end{array}
\end{equation}
递推只使用 Petersen 图的一维邻域表
\[
 n(k)=(0,3,5,6,6,8,9,9,10,10,10),\qquad0\le k\le10.
\]
附录给出递推的完整规范；因环面退化和平坦性，坐标下界推出任意
特征零子空间的下界。

令 \(p_d=\dim((E+H)^{(1)}/D)\)。由
\(J\subset(E+H)^{(1)}\)、\(D\cap J=S\) 及
\(\partial S\subset T\)，
\begin{equation}\label{n5:eq:p-lower}
 p_d\ge t+d-s\ge m_s+d-s.
\end{equation}
对同源同靶矩阵 \(A,B\)，
\[
 \rank(A-B)\ge \rank[A\ B]+\rank\!\begin{bmatrix}A\\B\end{bmatrix}
 -\rank A-\rank B.
\]
把它用于二次 Koszul 展平。列拼接秩为
\(2400+25d-p_d\)；行拼接秩为
\(25(100+h-s)-(100+d)\)。由于 \(J\subset H^{(1)}\)，这里只能且
只需使用 \(\rank K(C)\le24h\)。于是得到核心双商不等式
\begin{equation}\label{n5:eq:double-quotient}
 \rank K(R)\ge2400+(25d-p_d)-(25s-m_s).
\end{equation}

又由推论~\ref{n5:cor:intersection}，将 \(E\cap U\) 提升到
\(\bigoplus_iF_i\) 后投影到任意五块，其核至多一维，故
\begin{equation}\label{n5:eq:t51}
 t\le\dim(E\cap U)\le5\cdot10+1=51.
\end{equation}
式~\eqref{n5:eq:shadow-table} 的 \(m_{23}=52\) 因而给出 \(s\le22\)。
将式~\eqref{n5:eq:p-small}、\eqref{n5:eq:p-lower}、
\eqref{n5:eq:double-quotient} 逐格代入，若 \(s\le18\) 或 \(d\le8\)，
则要么延拓上下界冲突，要么
\(\rank K(R)>9\cdot240=2160\)。因此
\[
 19\le s\le22,\qquad d\ge9,\qquad h=t+d\le60.
\]
固定 \(s\) 时，允许的 \((t,d)\) 点构成边长
\(53-m_s\) 的整数三角形，故状态数为
\[
 \binom{53-m_s}{2}.
\]
由于 \(m_{19}=45\)、\(m_{20}=m_{21}=m_{22}=48\)，总数严格为
\[
 \binom82+3\binom52=28+30=58.
\]

\begin{lemma}[平方自由三次空间不含非零二元三次型]
\label{n5:lem:squarefree-not-binary}
设特征为零，\(z_1,\ldots,z_m\) 是 \(Y\) 的一组基，并令
\[
 J=\Span\{z_iz_jz_k:i,j,k\ \text{两两不同}\}\subset\Sym^3Y.
\]
若 \(L\subset Y\) 且 \(\dim L\le2\)，则
\[
 J\cap\Sym^3L=0.
\]
\end{lemma}

\begin{proof}
任意 \(f\in J\) 满足
\begin{equation}\label{n5:eq:squarefree-second-derivatives}
 \partial_{z_i^*}^2f=0\qquad(1\le i\le m).
\end{equation}
坐标协向量 \(z_i^*\) 在 \(L\) 上的限制张成 \(L^*\)。若
\(\dim L=1\)，写 \(f=c u^3\)，并选取一个在 \(L\) 上非零的坐标
协向量 \(\alpha\)。则
\(\partial_\alpha^2f=6c\alpha(u)^2u\)，故
式~\eqref{n5:eq:squarefree-second-derivatives} 给出 \(c=0\)。

若 \(\dim L=2\)，可选两个坐标协向量，其限制
\(\alpha,\beta\) 构成 \(L^*\) 的一组基。取 \(L\) 上的坐标 \(u,v\)，
使 \(\alpha=\partial_u\)、\(\beta=\partial_v\)，并写
\[
 f=a u^3+b u^2v+cuv^2+d v^3.
\]
式~\eqref{n5:eq:squarefree-second-derivatives} 对这两个坐标方向分别给出
\[
 0=\partial_u^2f=6au+2bv,
 \qquad
 0=\partial_v^2f=2cu+6dv.
\]
特征零立即推出 \(a=b=c=d=0\)。
\end{proof}

\begin{lemma}[近极大耦合]\label{n5:lem:coupling}
若上述 fixed-six 状态满足 \(h\ge57\)，则 \(H=U\)。
\end{lemma}

\begin{proof}
置
\[
 e=60-\sum_i\dim F_i,\quad
 k=\sum_i\dim F_i-\dim U,\quad c=\dim U-h.
\]
则 \(e+k+c=60-h\le3\)。反设 \(c>0\)。若 \(k=0\)，二次空间直和；
对任意三次关系逐项微分后落入二次直和，Euler 恒等式迫使每个分量
为零，故三次空间也直和。催化矩阵转置给 \(H=U\)，矛盾。

若 \(k=1\)，则 \(e\le1\)，所以每块二次维数至少九。转置给出非零
三次关系 \((q_i)\)，并且
\(\partial_vq_i=\lambda(v)r_i\)，其中 \((r_i)\) 生成唯一二次关系。
交换混合偏导再用 Euler 恒等式，得到 \(q_i\) 是同一线性式的纯立方。
但是二次维数至少九迫使因子张成至少四维；取四个独立因子为坐标，
第五个写成它们的线性组合，所有三因子积均不含任何坐标纯立方项，
故三次导数空间不含非零纯立方，矛盾。

若 \(k=2\)，必有 \(e=0\)。三次关系的导数元组若只张成一维，回到
上一段。若张成二维，存在同一二元空间 \(L=\langle\ell,m\rangle\)
使每个 \(q_i\in\Sym^3L\)。五因子独立时，其三次导数空间由因子坐标中
的平方自由三次单项式张成，引理~\ref{n5:lem:squarefree-not-binary}
迫使该分量 \(q_i=0\)。其余非零分量均来自四维因子空间
\(Y_i\)，且 \(F_i=\Sym^2Y_i\)。纯立方已被排除，故非零 \(q_i\) 的
本质空间恰为 \(L\subset Y_i\)。关系至少有两个非零分量，相应两个
\(F_i\) 共同包含三维 \(\Sym^2L\)，于是二次关系核至少三维，与
\(k=2\) 矛盾。最后 \(k\ge3\) 与 \(e+k+c\le3,c>0\) 直接矛盾。
\end{proof}

\subsection{58 态的完备路由}

表~\ref{n5:tab:routes} 给出互不重叠且穷尽的路由。其计数恒等式为
\[
38+1+1+9+2+1+1+5=58;
\]
每一类由式~\eqref{n5:eq:shadow-table}--\eqref{n5:eq:t51} 的整数
不等式直接定义，不读取历史汇总。

\begin{table}[ht]
\centering
\caption{fixed-six 的 58 态路由。}\label{n5:tab:routes}
\begin{tabular}{@{}lr@{}}
\toprule
排除机制 & 状态数\\ \midrule
\(s=19,d=9,t=45\) 的 Petersen--三角形等号计数 &1\\
全局 \(p_d\) 上界与式~\eqref{n5:eq:double-quotient} &38\\
\(d=9\) 等号几何 &1\\
\(p_{11}\le55\) &2\\
\(p_{12}\le80\) &1\\
\(p_{12}\le61\) &1\\
旗标零化子秩差（\(d=11,12\)） &5\\
\(d=10\) 的三轨道排除 &9\\ \midrule
合计 &58\\ \bottomrule
\end{tabular}
\end{table}

前 49 态不需要 d10 的有限轨道。其纯入口如下。Petersen 等号定理
给出 \(22\mapsto48\) 的唯一 row-star/column-star 旗标；删一三次基元
的 \(21\mapsto48\) 情形由父集计数恢复。若
\(S\subset D,T\subset E,\partial S\subset T\)，并对五平面 \(L\) 定义
\[
 W_L=\Span\{\partial_xs:s\in S, x\in L^\perp\},
\]
则在完整 \(S_{22},T_{48}\) 旗标上
\(\operatorname{codim}_T W_L\le4\)；在删点旗标上相应上界为七。
证明把假设失败的 \((L,Z)\) 放入射影闭关联簇，先退化旗标，再由
Borel 不动点定理把 \(L,Z\) 化为坐标子空间；五变量的行占用分拆
\(5,4+1,3+2,3+1+1,2+2+1,2+1+1+1,1^5\) 的父泛函最大数分别为
\(4,4,2,1,1,0,0\)。因此该上界作用于完整相对闭包，不存在
fixed-fibre 到 mixed--Rees 的跳步。

将六块关系核 \(K_F\) 投影到各 \(F_i\)，得到规范商映射
\[
 \rho:U\longrightarrow\bigoplus_iF_i/R_i,\qquad
 \dim\ker\rho\le5\dim K_F.
\]
链映射恒等式说明 \(\rho_i|_T\) 消灭 \(W_{L_i}\)。上述余维上界与
\(43>42\)、\(48>42\) 或相应删点不等式逐一排除 d11/d12 的五态。
这证明表~\ref{n5:tab:routes} 的 49 个非 d10 状态均不可能。

\subsection{九个 d10 状态与纯压缩分类}

余下状态为
\begin{equation}\label{n5:eq:d10states}
 s\in\{20,21,22\},\quad d=10,\quad
 t\in\{48,49,50\},\quad h=t+10.
\end{equation}
引理~\ref{n5:lem:coupling} 给 \(H=U\)。由
\(e+k=60-h\le2\)，至少四个 \(F_i\) 为十维。若一个十维块与 \(E\)
有一维交，则其九维商像位于 \(W=q(U)\) 中。定理
~\ref{n5:thm:one-intersection-flag} 给出 \(p(W)\le26\)，而式
~\eqref{n5:eq:p-lower} 在
\eqref{n5:eq:d10states} 中至少为 36，矛盾。因此这些十维块都横截
\(E\) 并同构映满同一个 \(W_{10}\)。若恰有四块，则关系亏损为零；
若有五块，避开唯一关系支撑中的一块；若有六块，从至多二维关系核中
选择两个坐标投影联合单射并删去它们。故总能选四个饱和块直和。

选 \(S\) 的二十维子空间，并在记录五平面 \(L\) 的闭关联簇上作
行列 Borel 退化。坐标极限可写成
\[
 \mathcal S\subset\binom{[5]}3\times\binom{[5]}3,\qquad
 \mathcal L\subset[5]\times[5],
\]
且
\begin{equation}\label{n5:eq:flagbounds}
 |\mathcal S|=20,\qquad |\partial\mathcal S|\le50,\qquad
 |\partial_{\mathcal L^\perp}\mathcal S|\le40,\qquad|\mathcal L|=5.
\end{equation}

\begin{lemma}[同时 elementary compression]\label{n5:lem:simshift}
对任意一对行指标或列指标作同一 elementary compression \(C\)，有
\[
 \partial_{(C\mathcal L)^\perp}(C\mathcal S)
 \subset C\bigl(\partial_{\mathcal L^\perp}\mathcal S\bigr).
\]
特别地，普通影子和可见影子的大小都不增加。
\end{lemma}

\begin{proof}
只证行压缩。把低、高指标记为 \(a,b\)，并在每个二点轨道中比较
影子。若被删行不属于 \(\{a,b\}\)，父点与子点具有同一轨道类型；
新低点来自至少一个旧轨道点，新高点保留则迫使两个旧父点同时存在。
若删去 \(a\)，新高子点只能来自同时含 \(a,b\) 的固定父点，而
\((a,c)\notin C\mathcal L\) 迫使旧 \(\mathcal L\) 的两个对应格点
都不遮蔽，故旧影子含高、低两点。若删去 \(b\)，新低子点只要求旧
两格点不同时遮蔽，故旧影子至少含轨道中的一点。固定子点的两种情形
相同。三种情形分别给出低点的“或”、高点的“且”和固定点的恒等
规则，恰为右端集合的 compression。列压缩由交换两坐标得到。
\end{proof}

反复压缩终止于真实的乘积分量序 ideal。令一维三子集 ideals 的最小
影子函数及其增量为
\[
\begin{gathered}
 n(k)=(0,3,5,6,6,8,9,9,10,10,10)_k,\\
 \Delta_j=n(j)-n(j-1)=(3,2,1,0,2,1,0,1,0,0)_j.
\end{gathered}                                             \tag{*}
\]
若终点的列纤维为 \(F_R\)、\(k_R=|F_R|\)，则
\[
 |\partial\mathcal S|\ge
 B(k):=\sum_{A\in\binom{[5]}2}
 n\!\left(\max_{R\supset A}k_R\right).                 \tag{**}
\]

下面的层缺陷恒等式把旧的 \(1405\) 轮廓递推整体替换。把 \(k_R\)
降序重排为分拆 \(\lambda=(\lambda_1,\ldots,\lambda_{10})\)，并置
\[
 H_j=\{R:k_R\ge j\},\qquad h_j=|H_j|=\lambda'_j,
 \qquad e(I)=|\partial I|-n(|I|).
\]
最小初段的加入次序为
\(012,013,023,123,014,024,124,034,134,234\)，新影子数为
\[
 c=(3,2,1,0,2,1,0,1,0,0),\qquad
 \Phi(\lambda)=\sum_{i=1}^{10}c_i n(\lambda_i).
\]

\begin{lemma}[层缺陷恒等式]\label{n5:lem:layerdefect}
对任意 order--reversing 大小轮廓 \(k\)，
\[
 \boxed{B(k)-\Phi(\lambda)
 =\sum_{j=1}^{10}\Delta_j e(H_j)\ge0.}                  \tag{LD}
\]
\end{lemma}

\begin{proof}
由 \(n(m)=\sum_{j\le m}\Delta_j\)，交换两个有限求和即得
\[
 B(k)=\sum_{j=1}^{10}\Delta_j|\partial H_j|.
\]
标准 Ferrers 族转置后的列高为 \(h_j\)，故同一层析给出
\(\Phi(\lambda)=\sum_j\Delta_jn(h_j)\)。两式相减即得 (LD)；
每个 \(H_j\) 是一维 ideal，故每个 \(e(H_j)\ge0\)。
\end{proof}

现在只需分类普通整数分拆。若 \(\lambda_1\ge5\) 且至少五行非零，
固定 \(a=\lambda_1\)，把上式的十一个 \(n(k)\) 逐项代入，所得最小值为
\[
\begin{array}{c|rrrrrr}
a&5&6&7&8&9&10\\ \hline
\min\Phi&51&54&51&54&54&54.
\end{array}                                                \tag{***}
\]
其完备性可由有限递推
\[
 Q_i(r,u)=\min_{0\le x\le\min(r,u)}
 \{c_i n(x)+Q_{i+1}(r-x,x)\}
\]
直接归纳：在 \(i=2,3,4,5\) 强制 \(x\ge1\)，取
\(Q_{11}(0,u)=0\)、\(Q_{11}(r,u)=+\infty\ (r\ne0)\)，再计算
\(3n(a)+Q_2(20-a,a)\)。因此 \(\Phi(\lambda)\le50\) 时，
\(\lambda\) 或其共轭至多四行；Ferrers 转置保证
\(\Phi(\lambda)=\Phi(\lambda')\)。写四行方向为 \((a,b,c,d)\)，则
\[
 \Phi(a,b,c,d)=3n(a)+2n(b)+n(c),\qquad a+b+c+d=20.
\]
按 \(a=5,\ldots,10\) 代入 (*)，全部解恰为
\begin{equation}\label{n5:eq:fourteen}
\begin{gathered}
 (10,10),\ (10,4,4,2),\ (10,4,3,3),\ (9,4,4,3),\\
 (8,4,4,4),\ (7,5,4,4),\ (5,5,5,5)
\end{gathered}
\end{equation}
及其七个共轭分拆。

还需证明这不是只对重排轮廓的比较。十点 poset 的唯一最小元为
\(012\)，故 \(F_{012}\supseteq F_R\) 对所有 \(R\) 成立；非空列数正是
\(|F_{012}|=\lambda_1\)。所以必要时转置后，实际族至多有四个非空行。
四行型的 \(H_1\) 是大小四的一维 ideal。此大小只有
\[
 J_4=(0,1,3,6),\qquad A_4=(0,1,2,3),\qquad e(A_4)=2.
\]
若 \(H_1=A_4\)，则 (LD) 的第一层已贡献
\(\Delta_1e(A_4)=6\)，而上述七型的 \(\Phi\ge48\)，矛盾。因此
\(H_1=J_4\)；包含于 \(J_4\) 的每个 ideal 按大小唯一，故所有
\(H_j=J_{h_j}\)。两行型的支撑 ideal 本来唯一。于是实际非空行就是链
\[
 012<013<023<123
\]
上的式~\eqref{n5:eq:fourteen}。

最后写相应嵌套纤维为
\(F_1\supseteq F_2\supseteq F_3\supseteq F_4\)。六个行二子集的最早
父行给出精确式
\[
 |\partial\mathcal S|
 =3|\partial F_1|+2|\partial F_2|+|\partial F_3|
 =\Phi(\lambda)+3e(F_1)+2e(F_2)+e(F_3).                 \tag{FD}
\]
十六个一维 ideals 中，非零缺陷只可能为
\[
\begin{array}{c|ccccc}
|I|&3&4&5&6&7\\ \hline
e(I)>0&1&2&1&1&1,
\end{array}
\]
且每个大小只有一个这种非标准 ideal；大小 \(0,1,2,8,9,10\) 唯一。
七型可用余量 \(50-\Phi\le2\)。故 (FD) 先强迫 \(F_1\) 标准；再强迫
大小四的 \(F_2\) 标准（代价至少四），而 \((7,5,4,4)\) 中大小五的
非标准 \(F_2\) 代价二也超过该型余量一。随后等大小嵌套强迫大多数
\(F_3,F_4\)；仅余的大小三非标准 ideal \((0,1,2)\) 不包含于
\(J_4=(0,1,3,6)\)，而 \((7,5,4,4)\) 的非标准大小四 \(F_3\) 代价二
仍超过余量一。大小二纤维唯一。故全部 \(F_i=J_{|F_i|}\)，
式~\eqref{n5:eq:fourteen} 确实分类原 shifted ideal 本身；旧
\(1405\) 轮廓遍历和七行纤维选择表不再承担证明责任。

下面不用七族乘七平面的四十九项表。对列二元组 \(B\) 置
\[
 P_k(B)=\{c\notin B:B\cup\{c\}\in J_k\},
\]
并以 \([u]=\{0,\ldots,u-1\}\)（\([0]=\varnothing\)）定义
\[
 g_{p,q}(u,v)=\#\{B:P_p(B)\cup P_q(B)\ne\varnothing,
                  \ P_p(B)\subseteq[u],\ P_q(B)\subseteq[v]\}.
\]
若四个非空行 \(012,013,023,123\) 的长度为
\(\lambda=(a,b,c,d)\)，五格 Ferrers 平面的行长为
\(\mu=(u_0,\ldots,u_4)\)，则六个行二元组逐一给出被平面完全杀掉的
父点数
\begin{equation}\label{n5:eq:six-edge}
\begin{aligned}
h_\lambda(\mu)={}&g_{a,b}(u_2,u_3)+g_{a,c}(u_1,u_3)
                  +g_{b,c}(u_1,u_2)\\
 &+g_{a,d}(u_0,u_3)+g_{b,d}(u_0,u_2)
                  +g_{c,d}(u_0,u_1).
\end{aligned}                                                   \tag{VE}
\end{equation}
确实，行二元组 \(01\) 的父行是 \(012,013\)，缺失指标为 \(2,3\)，
这给第一项；其余五项相同。并且
\[
 |\partial_{\mathcal L^\perp}\mathcal S|
 =|\partial\mathcal S|-h_\lambda(\mu),\qquad
 u_1\le2,\quad u_2,u_3\le1.                                 \tag{VA}
\]

令 \(\alpha_k(u)=g_{k,k}(u,u)\)。把十个 \(B\) 代入定义，得到四个
长度六的序列
\[
\begin{array}{c|rrrrrr}
u&0&1&2&3&4&5\\ \hline
\alpha_3&0&3&3&4&6&6\\
\alpha_4&0&0&1&3&6&6\\
\alpha_5&0&1&3&5&7&8\\
\alpha_{10}&0&0&0&1&4&10.
\end{array}                                                     \tag{VS}
\]
再记
\[
A=(0,0,0,1,4,4),\quad \rho=(0,0,1,2,3,3),\quad
\beta=(0,1,2,3,4,4),
\]
它们分别是式~\eqref{n5:eq:six-edge} 中出现的
\(g_{10,3},g_{4,3},g_{8,4}\) 截面。对四个非幸存族，式~(VE) 化为
\[
\begin{array}{c|l|c}
\lambda&h_\lambda(\mu)&\max h_\lambda\\ \hline
(10,4,3,3)&A(u_0)+\alpha_3(u_1)
 +\mathbf1_{u_2=1}\rho(u_0)+\mathbf1_{u_1=2,u_2=1}&7\\
(9,4,4,3)&2\mathbf1_{u_1=2}+\sigma(u_0,u_3)
 +(\mathbf1_{u_2=1}+\mathbf1_{u_1\ge1})\rho(u_0)&7\\
(8,4,4,4)&\beta(u_0)+\beta(u_1)+\beta(u_2)
 +\alpha_4(u_1)+2\alpha_4(u_2)&6\\
(5,5,5,5)&\alpha_5(u_1)+2\alpha_5(u_2)+3\alpha_5(u_3)&6,
\end{array}                                                     \tag{VB}
\]
其中
\[
 \sigma(u,0)=(0,0,2,2,4,4)_u,\qquad
 \sigma(u,1)=(0,0,3)_u\quad(u\le2).
\]
这里的上界不含枚举：若 \(u_2=1\)，按 \(u_1=1,2\)；若
\(u_2=0\)，按 \(u_1=0,1,2\) 分开，面积恒等式
\(\sum u_i=5\) 直接代入 (VS)。第三、四行只需再分
\(u_1\le1\) 与 \(u_1=2\)。因此四族的可见影子依次至少为
\(41,41,42,42\)。

余下三族的等号也由父集而非四十九项表确定。
类型 \(0=(10,10)\) 有
\(h_0=2\alpha_{10}(u_0)\)，故阈值四十强迫 \(u_0=5\)。
类型 \(1=(10,4,4,2)\) 的行 \(0\) 内只有一个重复二点父集
\(\{00,01\}\) 和五个三点父集；四格至多包含其中四个三点父集，
等号唯一强迫四格行为 \(\{00,01,02,03\}\)。再杀两个跨行父集时，
第五格恰为 \(10,11,20,21\)，所以最大杀点数为八。
类型 \(13=(7,5,4,4)\) 中，达到九点只可能使用变量行 \(0,1\)；若
三格、二格的列集分别为 \(A,B\)，则
\[
\begin{aligned}
h={}&|A\cap\{0,1\}|+|B\cap\{0,1\}|
 +\binom{|A\cap\{0,1,2\}|}{2}\\
&+\binom{|B\cap\{0,1,2\}|}{2}
 +\mathbf1_{B\subset A\cap\{0,1,2,3\}}\le9.
\end{aligned}
\]
且等号恰为 \((A,B)=(012,01)\) 及交换两变量行所得者。
这些父集的完整短证见式~(F5)--(F6)；旧四十九项表仅作整数交叉诊断。
结合普通影子阈值，shifted 好旗标在转置下只剩
\begin{equation}\label{n5:eq:survivors}
\begin{array}{c|c|c|c}
\text{类型}&\mathcal S&\mathcal L&
(|\partial\mathcal S|,|\partial_{\mathcal L^\perp}\mathcal S|)\\ \hline
0&(10,10)&(5)&(50,30)\\
1&(10,4,4,2)&(4,1)&(48,40)\\
13&(7,5,4,4)&(3,2)&(49,40).
\end{array}
\end{equation}

最后反向读取原 elementary-shift 序列。对一步逆压缩，每个可移动
二轨道给一个取向变量。若两个取向相反时某个目标外二次点 \(w\)
必出现，就以 \(w\) 连接二变量；\(w\) 只有至多九个父点，故每条边
是至多九项二文字析取的直接检查。

类型 \(0\) 不需要见证表。以补集把第二坐标三元组识别为 Petersen 图
\(P=KG(5,2)\) 的顶点；若十变量移动块的高取向集合为
\(\varnothing\ne X\ne V(P)\)，则
\[
N_P(X)\cap N_P(X^c)\ne\varnothing,
\]
否则长二路端点同色，而 Petersen 图的连通性和五圈会迫使 \(X\) 为空
或全体。因此每个十变量块统一；两个块同时移动而统一位相反时，第一
坐标的一维影子从五增到六，故全影子为六十。归一化 \(S=S_0\) 后，
若五格 \(L\) 分布在第零行与第 \(b\) 行，后者有 \(k\) 格，则可见
影子为
\[
\begin{cases}
50-2\binom{k}{3}-2\binom{5-k}{3},&b=1,\\
50-2\binom{5-k}{3},&b=2,3,4.
\end{cases}
\]
阈值四十只留下同轨的完整行，故 orbit--0 的原九十二条森林边全部
删除。

类型 \(1,13\) 也不再需要见证表。仍以补集识别 Petersen 图，并把
第一坐标纤维写成 \(S_R\)。对每个二元组 \(A\)，有恒等式与两级下界
\[
(\partial S)_A=N_P(U_A),\qquad
U_A=\bigcup_{R\sim A}S_R,\qquad
|\partial S|\ge\sum_A n(|U_A|)\ge
\sum_A n\!\left(\max_{R\sim A}|S_R|\right),                 \tag{F1}
\]
其中
\[
n(k)=(0,3,5,6,6,8,9,9,10,10,10)_k.                         \tag{F2}
\]
一个 ground--set 反射在 Petersen 顶点上只有三个二轨道；若目标低、
高纤维为 \(J_\alpha\supset J_\beta\)，逆像只把某个整数 \(x\) 个点
从低纤维移到高纤维。因此式~(F1) 只含至多三个整数。这里不再逐方向
代入。对任意大小轮廓 \(k\)，置
\[
H_j=\{R:k_R\ge j\},\qquad
\eta(X)=|N_P(X)|-n(|X|),
\]
并把 \(k\) 降序重排为 \(\lambda\)。层析证明完全不使用 shifted
假设，故层缺陷恒等式加强为
\[
B(k)-\Phi(\lambda)
=\sum_{j=1}^{10}\Delta_j\eta(H_j)\ge0.                         \tag{F3}
\]
所以 \(B(k)\le50\) 先强迫 \(\lambda\) 属于
式~\eqref{n5:eq:fourteen} 的七个分拆或其共轭。

按 Petersen 边次序
\(01,02,03,04,12,13,14,23,24,34\)，四个行、列大小轮廓为
\[
\begin{aligned}
r_1&=(0,0,0,2,0,0,4,0,4,10),&
c_1&=(1,1,1,3,1,1,3,1,4,4),\\
r_{13}&=(0,0,0,4,0,0,4,0,5,7),&
c_{13}&=(0,0,1,4,0,1,4,2,4,4).
\end{aligned}
\]
Petersen 图中四点邻域最小集恰为五个四点星；在一个换位的三个二轨道
上只反射非空真子集时，邻域从六增到八。七点邻域最小集恰为
\(N_P(v)^c\)。把上述四个轮廓的至多三个值对
\[
(\alpha,\beta)\longmapsto(\alpha-x,\beta+x)
\]
与十四个低分拆比较，再用这两个等号分类，所有非统一低 \(B\) 情形
恰压成三种移动模式：
\[
\begin{array}{c|c|c|c}
\text{轮廓}&\text{值对}&B\le50\text{ 的计数}&\text{方向}\\ \hline
r_1&(4,2)&0,1,2&r01,r02\\
r_1&(10,2)&0,1,2,3,5,6,7,8&r03\\
c_{13}&(2,0),(1,0)&(0,0),(1,0),(1,1),(2,1)&c03,c13.
\end{array}
\]
涉及四点星中心的混合端点由式~(F3) 至少增加
\(10,12,6,2\)，其余单块或七点混合情形也超过五十。因此旧四十二个
方向记录不再承担证明责任。这里 \(rab\) 表示 row
\(a\leftarrow b\)，\(cab\) 同理。

前两方向在 \(J_4\setminus J_2\) 中各有两个自由点，故各留下两个
非统一状态。对 \(r03\)，若
\(X\subset J_{10}\setminus J_2\) 且
\(b=|X\setminus J_4|\)，式~(F1) 的中项成为
\[
16+2\{n(10-b)+n(4+b)\}
=(48,52,54,52,54,52,48)_b.                                  \tag{F4}
\]
故只剩 \(b=0,6\)，共六个非统一族取向。对 \(c03,c13\)，三个取向位
\(z_0,z_1,z_2\) 满足：\(z_0\ne z_2\) 时中项至少五十二；
\(000,010,101,111\) 时分别为 \(49,50,50,49\)。故普通影子界总共只
留下十四个非统一族取向。

固定统一族后，令
\[
M_q(S)=\{(r,c):(A\cup\{r\},B\cup\{c\})\in S\},\qquad q=(A,B).
\]
若 \(h_S(L)\) 是满足
\(\varnothing\ne M_q(S)\subset L\) 的 \(q\) 数，则
\[
|\partial_{L^\perp}S|=|\partial S|-h_S(L).                    \tag{F5}
\]
先在每行压缩 \(L\)，再排序嵌套行；\(S\) 已 shifted，故可见影子不增。
按 \(L\) 的行分拆，式~(F5) 给出
\[
\begin{array}{c|rrrrrrr}
&5&4+1&3+2&3+1+1&2+2+1&2+1+1+1&1^5\\ \hline
\max h_{S_1}&7&8&5&5&2&2&0\\
\max h_{S_{13}}&5&6&9&6&7&4&2.
\end{array}                                                     \tag{F6}
\]
等号父集直接给出：类型 \(1\) 的四个好平面恰为
\[
\{00,01,02,03\}\cup\{sb\},\qquad s\in\{1,2\},\ b\in\{0,1\},
\]
类型 \(13\) 的两个好平面恰为
\[
\{00,01,02,10,11\},\qquad \{00,01,10,11,12\}.
\]
各组由相应 \(S\) 的稳定子传递。最后把式~(F5) 用于上述十四个兼容
逆像时不再调用十四状态表。置 \(y_i=1-z_i\)。对类型 \(1\) 的方向
\(r01\)，两个非统一族取向互为全互补，其中一个的被杀父集数为
\[
h_{01}=e_2(z)+z_0z_1z_2+e_2(y)+y_0y_1y_2+y_0+y_1+y_2.
                                                                    \tag{F7a}
\]
按 \(w=z_0+z_1+z_2=0,1,2,3\)，右端为 \(7,3,2,4\)。对方向
\(r02\)，一个取向的四变量公式为
\[
h_{02}=e_3(z_0,z_1,z_2,z_3)+e_3(y_0,y_1,y_2,y_3)
 +z_0y_1+y_0(y_1+y_2+y_3).                                  \tag{F7b}
\]
若 \(z_0=0\)，令 \(w=z_1+z_2+z_3\)，右端是
\(\binom w3+\binom{4-w}3+3-w\le7\)；若 \(z_0=1\)，它是
\(\binom{1+w}3+\binom{3-w}3+1-z_1\le4\)。互补取向结论相同。

对 \(r03\) 的六个状态，记
\(K=(J_{10}\setminus J_2)\setminus J_4\)。以 \(ab,cd\) 表示二次点
\((\{a,b\},\{c,d\})\)，并置
\[
Q_0=\{12,04;12,14;12,24;12,34\}.
\]
当兼容平面遍历时，可能被杀的二次点全体恰为
\[
\begin{array}{c|c}
X&\mathcal Q_X\\ \hline
\{3\}&Q_0\cup\{23,12;23,13;23,23\}\\
K\cup\{6\}&Q_0\cup\{02,12;02,13;02,23\}\\
\{6\},\{3,6\},K,K\cup\{3\}&Q_0.
\end{array}                                                     \tag{F7c}
\]
故任一固定平面至多杀七点。最后，类型 \(13\) 的 \(c03,c13\) 有同一
二变量公式
\[
h_{13}=7-2z_0-z_1+2z_0z_1\le7,                              \tag{F7d}
\]
另一个非统一取向仍由全互补得到。由于相应普通影子大小分别为
\(48,50\)，五个方向统一得到
\[
|\partial_{L^\perp}S_1|\ge48-7=41,\qquad
|\partial_{L^\perp}S_{13}|\ge50-7=43>40.                    \tag{F7}
\]
因此每一步逆像仍在同一行列置换及转置轨道，
式~\eqref{n5:eq:survivors} 分类所有满足~\eqref{n5:eq:flagbounds}
的旗标，而非仅分类其 shifted 极限。式~(F1)--(F4) 的层缺陷与移动
模式展开见内嵌文件
\path{n5_inverse_shift_layer_orbit_reduction_20260810.md}；
式~(F5)--(F6) 的统一父集展开见
\path{n5_flag_fibre_profile_structural_reduction_20260810.md}，十四状态的
式~(F7a)--(F7d) 展开见
\path{n5_exceptional_visible_parent_reduction_20260810.md}；旧的
\(250\) 条 witness 边仅保留为交叉诊断，不再是证明依赖。

\subsubsection{orbit--1：四个临界商终端}

令 \(p_{I,J}=\perm(X_{I,J})\)，并置
\[
\mathcal J_{10}=\binom{[5]}3,\qquad
\mathcal J_4=\{012,013,023,123\},\qquad
\mathcal J_2=\{012,013\}.
\]
\[
\begin{aligned}
 S_0={}&\Span\{p_{012,J}:J\in\mathcal J_{10}\}\\
 &+\Span\{p_{013,J},p_{023,J}:J\in\mathcal J_4\}\\
 &+\Span\{p_{123,J}:J\in\mathcal J_2\},
 \qquad T_0=\partial S_0 .
\end{aligned}
\]
这就是类型 \(1=(10,4,4,2)\) 的坐标旗标。记 \(q:\Sym^2V\to Q\)
为商映射，并写
\[
\begin{aligned}
S_{ia}&=q(x_{ia}^2),&
R_{i;ab}&=q(x_{ia}x_{ib}),\\
C_{ij;a}&=q(x_{ia}x_{ja}),&
X_{ij;ab}&=q(x_{ia}x_{jb})=-q(x_{ib}x_{ja})
\end{aligned}
\]
（后式取 \(i<j,a<b\)）。令
\[
 M=\langle x_{00},x_{01},x_{02},x_{03}\rangle,\qquad
 L^\dagger=M+\langle x_{10}\rangle .
\]
因 \(E\cap\Sym^2L^\dagger=0\)，\(q(\Sym^2L^\dagger)\) 是十五条
互异的一维商权。对其中任一十元子集 \(W\)，仍写
\(p(W)=\dim(E+W)^{(1)}-\dim D\)。下面给出不使用十元集穷举的纯图
分类。

\begin{proposition}[orbit--1 十五权终端图定理]
\label{n5:prop:orbit1terminal}
十五权宇宙的全部 \(\binom{15}{10}=3003\) 个十元子集满足
\(p(W)\le36\)。等号恰有四个：
\[
 W_M=q(\Sym^2M)
\]
以及三个在固定列 \(0\) 后互为列置换的平面
\[
\begin{aligned}
W_b=\Span(&\{R_{0;ac}:0\le a<c\le3\}\cup\{S_{00},C_{01;0}\}\cup\\
&\{S_{0a}:a\in\{1,2,3\}\setminus\{b\}\}),\qquad b=1,2,3.
\end{aligned}
\]
\end{proposition}

\begin{proof}
写
\[
y_i=x_{0i}\quad(0\le i\le3),\qquad z=x_{10}.
\]
把所选平方权记为
\(S\subset\{0,1,2,3,z\}\)，把所选非平方权记为同一五顶点集上的
简单图 \(G\)；于是
\[
 |S|+|E(G)|=10.                                      \tag{T1}
\]
对 \(1\le i<j\le3\)，记
\[
\begin{aligned}
a_i&=\mathbf1_{\{0i\in E(G)\}},&
b_{ij}&=\mathbf1_{\{ij\in E(G)\}},\\
c&=\mathbf1_{\{0z\in E(G)\}},&
x_i&=\mathbf1_{\{iz\in E(G)\}}.
\end{aligned}
\]
令 \(H=G[\{0,1,2,3\}]\)，令 \(\tau(H)\) 为 \(H\) 的三角形数，
\(\tau_z(G)\) 为含 \(z\) 的三角形数，并置
\[
\begin{aligned}
q_4(G)&=\sum_{i<j}a_ia_jx_ix_j(1-b_{ij}),\\
\eta(G)&=\sum_i x_i+c\sum_i a_i(1-x_i).
\end{aligned}                                                   \tag{T2}
\]

先证明闭式
\[
\begin{aligned}
p(W)={}&5\tau(H)+\tau_z(G)+q_4(G)\\
&+\epsilon_0\bigl(1+d_G(0)+\eta(G)\bigr)
 +\sum_{v\in\{1,2,3,z\}}\epsilon_v\bigl(1+d_G(v)\bigr),
\end{aligned}                                                   \tag{T3}
\]
其中 \(\epsilon_v=\mathbf1_{\{v\in S\}}\)。

固定一个三次行列环面权，以该权的三次单项式为顶点。divided-power
偏导再模去 \(E\) 后，未选二次权给零锚或带符号边；选中该权就删除
相应锚或边。局部相对核是删除后无锚且符号相容的连通分量数。按列重数
\(3,2+1,1+1+1\) 分类并收缩标签不在这十五权中的固定边。为把
“穷尽”也写成可手检的有限表，先按行、列重数记录非零局部块；表项
括号内是其中的非单点约化图数：
\[
\begin{array}{c|ccc|c}
 &3&2+1&1+1+1&\text{行和}\\ \hline
3&5&12&4&21\\
2+1&2&6+(3)&13+(3)&27\\ \hline
\text{列和}&7&21&20&48
\end{array}                                                     \tag{T4a}
\]
这里第一行依次是五个平方立方块、十二个同一基底行内的平方--边块、
四个基底三角形块；第二行第一格是边 \(0z\) 的两个平方--边块，
第二格是六个普通平方--边块与三个 (T5) 块，第三格是每个基底三角形
的四个逐行换位块，其中三个是 (T6) 块。若行重数为
\(1+1+1\)，每个分量仍含一个标签不在十五权中的零锚；上表其余
格子的未列分量也保留这种零锚。因此它们的相对核恒为零。这同时
证明 (T4a) 穷尽了全部可能产生非零相对核的三次环面权。

把这四十八个块的局部核函数作 Möbius 展开，全部非零项只有
\[
\begin{array}{c|c|c}
\text{项族}&\text{个数}&\text{系数}\\ \hline
s_v&5&1\\
s_ue_{uv}\ (u\ne v)&20&1\\
s_0x_i&3&1\\
s_0a_ic&3&1\\
s_0a_icx_i&3&-1\\
\prod_{e\subset T}e,\quad T\subset\{0,1,2,3\},\ |T|=3&4&5\\
ca_ix_i&3&1\\
b_{ij}x_ix_j&3&1\\
a_ia_jx_ix_j&3&1\\
a_ia_jb_{ij}x_ix_j&3&-1.
\end{array}                                                     \tag{T4}
\]
这里 \(s_v,e_{uv}\) 是相应指示量。表中只有两个非单点约化图。对每个
\(i\)，两顶点
\[
 \{y_0^2x_{1i},\,y_0y_i z\}
\]
由标签 \(x_i\) 相连，且分别被 \(s_0\) 与 \(c,a_i\) 锚定；逐个
删标所得 Möbius 多项式为
\[
 s_0x_i+s_0a_ic+a_icx_i-s_0a_icx_i.                \tag{T5}
\]
对 \(i<j\)，三顶点
\[
 \{y_0y_ix_{1j},\,y_0y_jx_{1i},\,y_iy_jz\}
\]
中，\(x_j,x_i\) 分别连接第一、第三点及第二、第三点，前两点还有一条
固定边，而 \(a_i,a_j,b_{ij}\) 分别锚定三点；故其多项式为
\[
 a_ia_jb_{ij}+b_{ij}x_ix_j+a_ia_jx_ix_j
 -a_ia_jb_{ij}x_ix_j.                              \tag{T6}
\]
平方立方块给 \(s_v\)，平方--边块给 \(s_ue_{uv}\)。每个基底三角形
另有一个全在第零行的单项式块和四个逐行换位块，各贡献三边乘积，
所以系数为五。(T5)--(T6) 与 (T4a) 中的单点约化图逐格对应，故
(T4) 的五十个非零 Möbius 项完备；按平方、三角形和四圈分组即得
(T3)。

现在置
\[
 m=|E(G)|,\qquad s=|S|=10-m;
\]
故 \(5\le m\le10\)。若 \(H\ne K_4\)，则
\[
 \tau(H)\le2,\qquad
 \tau_z(G)+q_4(G)
 =\sum_i ca_ix_i+
 \sum_{i<j}x_ix_j\bigl(b_{ij}+(1-b_{ij})a_ia_j\bigr)\le6.
                                                               \tag{T7}
\]
所以纯边部分至多十六。记
\[
 w_0=1+d_G(0)+\eta(G),\qquad w_v=1+d_G(v)\quad(v\ne0).
\]
则
\[
0\le\eta\le3,\quad 1\le w_0\le8,\quad
1\le w_v\le5\ (v\ne0),\quad
\sum_vw_v=5+2m+\eta.                                  \tag{T8}
\]
按 \(m=5,6,7,8,9\)，最大的 \(s\) 个平方权之和及相应 \(p\)-上界为
\[
\begin{array}{c|rrrrr}
m&5&6&7&8&9\\ \hline
s&5&4&3&2&1\\
\text{平方和上界}&18&19&18&13&8\\
p(W)\text{ 上界}&34&35&34&29&24.
\end{array}                                                     \tag{T9}
\]
\(m=6\) 一列由总和删去至少为一的最小权得到；\(m=10\) 时
\(|E(H)|\le5\) 不可能。因此等号必须有 \(H=K_4\)。

最后设 \(H=K_4\)，并令
\[
 \rho=x_1+x_2+x_3,\qquad r=c+\rho=m-6.
\]
此时
\[
\tau(H)=4,\qquad \tau_z=c\rho+\binom\rho2,\qquad q_4=0,
\]
而平方权为
\[
\begin{array}{c|ccc}
&w_0&w_i\ (1\le i\le3)&w_z\\ \hline
c=0&4+\rho&4+x_i&1+\rho\\
c=1&8&4+x_i&2+\rho.
\end{array}                                                     \tag{T10}
\]
选取最大的 \(s=4-r\) 个权，只有八行：
\[
\begin{array}{c|c|c|c|c|c}
r&c&\rho&s&\text{纯边部分}&\text{平方和}\quad p_{\max}\\ \hline
0&0&0&4&20&16\quad36\\
1&0&1&3&20&14\quad34\\
1&1&0&3&20&16\quad36\\
2&0&2&2&21&11\quad32\\
2&1&1&2&21&13\quad34\\
3&0&3&1&23&7\quad30\\
3&1&2&1&23&8\quad31\\
4&1&3&0&26&0\quad26.
\end{array}                                                     \tag{T11}
\]
第一行强迫选择 \(K_4\) 六边和四个基底平方，正是 \(W_M\)。第三行
强迫再选 \(c=C_{01;0}\)，平方为 \(s_0\) 与 \(s_1,s_2,s_3\) 中任意
两个，恰是三个 \(W_b\)。其余行严格小于三十六，命题得证。
\end{proof}

这里补明图定理与任意 orbit--1 分支之间的衔接。route~4 的等号链给出
共同十维商像 \(W\) 及 \(p(W)\ge36\)。沿产生上述坐标旗标的行列环面
平坦退化后，核维的上半连续性给 \(p(W^*)\ge p(W)\)。同时
\(E\cap\Sym^2L^\dagger=0\)，所以四个饱和块在特殊纤维均不发生商秩
下降，且仍同构映满同一个坐标十平面 \(W^*\)。命题
~\ref{n5:prop:orbit1terminal} 于是强迫
\(p(W^*)=p(W)=36\)，并且 \(W^*\) 必为下述四个终端之一。这正是从
非坐标反例到有限终端表的完备桥梁。

三个 \(W_b\) 由同一个长度二局部论证处理。取其中
\[
\begin{aligned}
W_0=\Span\{&
S_{00},S_{01},S_{02},
R_{0;01},R_{0;02},R_{0;03},\\
&R_{0;12},R_{0;13},R_{0;23},C_{01;0}\}.
\end{aligned}                                                     \tag{O1}
\]
其中三个 square 权 \(S_{00},S_{01},S_{02}\)、同行权
\(R_{0;03}\) 与同列权 \(C_{01;0}\) 依次强迫
\[
 x_{00},x_{01},x_{02},x_{03},x_{10}\in L.
\]
这些商权空间均为一维，故四个最大五平面见证在特殊纤维共同落到唯一
\[
 L_0=\langle x_{00},x_{01},x_{02},x_{03},x_{10}\rangle.
\]
又 \(\dim W_{L_0}^{\perp}=8\)，所以关联胚中的八平面也唯一：
\(Z_0=W_{L_0}^{\perp}\)。在相应图表中写
\[
 L=M+\langle x_{10}+Bx_{11}+Cx_{20}+Dx_{21}\rangle,
 \qquad M=\langle x_{00},x_{01},x_{02},x_{03}\rangle.
\]
线性化零化子关联的每条方程只有一个系数或两个系数之差。具体地，
以 \([I;J]\) 表示行、列二子集给出的 \(T_0\) 坐标；若
\(p=p_{K,L}\in S_0\) 且 \(I\subset K,J\subset L\)，记
\[
m_p([I;J])=x_{K\setminus I,\ L\setminus J}.
\]
用 \(L(u\to v)\) 和 \(Z(\alpha\to\beta)\) 表示两个 Grassmann 图
坐标。对 \(\alpha\in Z_0,p\in S_0,v\notin L_0\)，线性化方程正是
\[
\mathbf1_{m_p(\alpha)\ {\rm defined}}L(m_p(\alpha)\to v)
-\!\!\sum_{\substack{\beta\in T_0\setminus Z_0\\m_p(\beta)=v}}
 Z(\alpha\to\beta)=0.                                      \tag{O2}
\]
右侧和至多一项，所以式~(O2) 是单项锚或二项边。逐格使用这一同一
判据，四百二十个顶点中没有直接锚的二十三个恰为
\[
\begin{array}{c|l}
\text{类型与源}&\text{目标}\\ \hline
L(x_{02}\to\cdot)&x_{22}\\
L(x_{03}\to\cdot)&x_{23}\\
L(x_{10}\to\cdot)&x_{11},x_{20},x_{21},x_{30},x_{31}\\
Z([12;04]\to\cdot)&[01;04],[02;04]\\
Z([12;14]\to\cdot)&[01;14],[02;14]\\
Z([12;24]\to\cdot)&[01;24],[02;24],[12;23]\\
Z([12;34]\to\cdot)&[01;34],[02;34],[12;23]\\
Z([23;12]\to\cdot)&[13;02],[13;12],[23;02]\\
Z([23;13]\to\cdot)&[13;03],[13;13],[23;03].
\end{array}                                                  \tag{O3}
\]
其中十四个又分别与下列方括号中的显式锚位于同一分量：
\[
\begin{aligned}
[L(x_{00}\to x_{20})]&\sim
\{L(x_{02}\to x_{22}),L(x_{03}\to x_{23}),\\[-1mm]
&\hspace{8mm} Z([12;04]\to[01;04]),Z([12;14]\to[01;14]),\\[-1mm]
&\hspace{8mm} Z([12;24]\to[01;24]),Z([12;34]\to[01;34])\},\\
[L(x_{01}\to x_{11})]&\sim
\{Z([12;04]\to[02;04]),Z([12;14]\to[02;14]),\\[-1mm]
&\hspace{8mm} Z([12;24]\to[02;24]),Z([12;34]\to[02;34])\},\\
[L(x_{02}\to x_{04})]&\sim\{Z([12;34]\to[12;23])\},\\
[L(x_{03}\to x_{04})]&\sim\{Z([12;24]\to[12;23])\},\\
[Z([23;12]\to[12;12])]&\sim\{L(x_{10}\to x_{30})\},\\
[Z([23;12]\to[12;02])]&\sim\{L(x_{10}\to x_{31})\}.
\end{aligned}                                                 \tag{O4}
\]
唯一三个无锚分量为
\[
\begin{aligned}
&\{L(x_{10}\to x_{11}),Z([23;12]\to[23;02]),
                         Z([23;13]\to[23;03])\},\\
&\{L(x_{10}\to x_{20}),Z([23;12]\to[13;12]),
                         Z([23;13]\to[13;13])\},\\
&\{L(x_{10}\to x_{21}),Z([23;12]\to[13;02]),
                         Z([23;13]\to[13;03])\}.
\end{aligned}                                                 \tag{O5}
\]
式~(O2) 逐项产生式~(O3)--(O5)，故这是一张可手检的整数关联表，而非
矩阵秩断言。切空间维数恰为三，而上述
\(\mathbf P^3\) 运动显式包含在关联胚中；正则局部环中的同维闭浸入
遂说明固定旗标的形式局部胚就是该 \(\mathbf P^3\) 胚。

加入共同商条件 \(W_0\subset q(\Sym^2L)\) 后不需要 Gröbner 基。
式~(O1) 的前九个生成元已在 \(q(\Sym^2M)\) 中，只需提升
\(C_{01;0}=q(x_{00}x_{10})\)。写一般提升为
\[
u=u_M+a_0x_{00}y+a_1x_{01}y+a_2x_{02}y+a_3x_{03}y+a_4y^2,
\qquad u_M\in\Sym^2M .
\]
在 \(Q\) 中依次比较
\[
C_{01;0},\ X_{01;01},\ C_{02;0},\ X_{01;02},\
X_{01;03},\ S_{10}
\]
六个互异权，得到
\[
a_0=1,\qquad a_1=B,\qquad C=0,\qquad
a_2=a_3=a_4=0.
\]
再比较 \(C_{01;1}\) 与 \(X_{02;01}\) 得
\[
B^2=0,\qquad D-BC=0,
\]
故 \(D=0\)。反过来，
\[
u=x_{00}y+Bx_{01}y
\]
模 \((B^2,C,D)\) 的商像正是 \(C_{01;0}\)。因此这不是根集计算，而是
完整的方案论系数比较，消元理想恰为
\begin{equation}\label{n5:eq:length2}
 (B^2,C,D),\qquad
 \mathcal O_{\mathrm{fiber},L_0}\simeq\Q[B]/(B^2).
\end{equation}
把 \((S,T,W,L,Z)\) 全部记录在射影 Grassmann 乘积的闭关联簇中。
设基与总空间在特殊点的完成局部环分别为 \(R,A\)。由
式~\eqref{n5:eq:length2}，\(A/\mathfrak m_RA\) 的长度为二。提升其
两个基元素并逐次模 \(\mathfrak m_R^jA\) 表示，完备性与拓扑 Nakayama
给出 \(A\) 是由两个元素生成的有限 \(R\)-模。故任意 DVR 基变换后，
几何泛纤维的长度、从而互异几何点数至多为二。这一完成局部环论证
允许基参数与垂直参数的乘积项。
给定 \((L,W)\) 后，\(q|_{\Sym^2L}\) 在邻域内单射，故每点至多给一个
\(F\)。四个直和饱和块必须互异，矛盾。该论证同时覆盖 orbit--1 在
\(t=48,49,50\) 的后代，且正面处理 base--vertical 混合项。

\subsubsection{orbit--1 的共同包络终端 \(W_M\)}

命题~\ref{n5:prop:orbit1terminal} 的第四个等号终端是
\[
 W_M=q(\Sym^2M),\qquad
 N=\langle x_{10},x_{11},x_{20},x_{21}\rangle .
\]
固定旗标的五平面特殊纤维为
\(L^*=M+\langle y\rangle\)，其中 \([y]\in\mathbf P(N)\)。下面的
赋值引理排除此前长度二论证没有覆盖的共同包络。

\begin{lemma}[同行全阶碰撞]
\label{n5:lem:samerow}
令 \(\mathcal O=k[[t]]\)、\(K=k((t))\)。设
\[
\begin{aligned}
 F_i&\subset\Sym^2L_i,& F_i\cap E_K&=0,\\
 q(F_0)&=q(F_1)=W,& \dim W&=10,\quad i=0,1,
\end{aligned}
\]
且两个特殊二次空间均为 \(\Sym^2M\)，两个特殊五平面均为
\(M\) 加上同一外部行中的一条坐标线；两条线允许相同。则
\[
 \dim_K(F_0\cap F_1)\ge9.
\]
\end{lemma}

\begin{proof}
置
\[
 d=\dim(L_0\cap L_1),\qquad s=\dim(F_0\cap F_1).
\]
五平面交给
\[
 s\le\binom{d+1}{2},\qquad
 \dim\bigl(E_K\cap(F_0+F_1)\bigr)=10-s.               \tag{SR1}
\]
若 \(d=5\)，共同商像与
\(\dim(E\cap\Sym^2L_0)\le1\) 已直接给 \(s\ge9\)。以下设
\(0\le d\le4\)。

先记录所需的坐标矩形界。若格点集 \(S\) 含
\(M=\{00,01,02,03\}\)，令 \(\epsilon\) 记录 \(04\) 是否属于
\(S\)，并令 \(n_i\) 为第 \(i>0\) 行的格点数，则
\[
\begin{aligned}
R(S)&=\sum_{i>0}\binom{|S_0\cap S_i|}{2}
 +\sum_{0<i<j}\binom{|S_i\cap S_j|}{2}\\
&\le\sum_{i>0}\binom{\min(n_i,4+\epsilon)}2
 +\sum_{0<i<j}\binom{\min(n_i,n_j)}2 .
\end{aligned}                                                   \tag{SR2}
\]
这里 \(R(S)=\dim(E\cap\Sym^2\langle S\rangle)\)。对新增格点总数
\(q=\epsilon+\sum n_i\le6\)，只需对 \(q-\epsilon\) 的整数分拆
代入 (SR2)，得到
\[
\begin{array}{c|rrrrrrr}
q&0&1&2&3&4&5&6\\ \hline
\max R&0&0&1&3&6&6&10.
\end{array}                                                   \tag{SR3}
\]
最后一个等号唯一来自 \(\epsilon=1\) 且某一外部行含五格，即完整
两行块。闭条件 \(\dim(E\cap\Sym^2K)\ge h\) 对行列环面不变；若其
轨道闭包只有上述一个固定点，则正维完备 toric 轨道闭包至少有两个
固定点的标准事实迫使原点本身固定。故 (SR3) 对一般六至十平面也可
通过环面特殊化使用，而不只是坐标诊断。

由 (SR1)--(SR3)，\(d=1,2,3\) 分别要求至少 \(9,7,4\) 个矩形，而
相应上界只有 \(6,6,3\)，均不可能。若 \(d=0\)，所有不等式取等，
饱和和空间的特殊纤维必为
\[
 K_*=\langle r_0,r_1\rangle\otimes C,\qquad \dim C=5,
\]
且其十维 \(E\)-交专门化为
\((r_0\odot r_1)\otimes C_{\rm off}\)，其中
\(C_{\rm off}\subset\Sym^2C\) 是对称零对角矩阵空间。

固定坐标补空间后，Nakayama 引理把泛纤维和空间写成图
\[
 K=\Gamma(h),\qquad
 h(r_0\otimes c)=\sum_{u=2}^4r_u\otimes P_uc,\quad
 h(r_1\otimes c)=\sum_{u=2}^4r_u\otimes Q_uc.          \tag{SR4}
\]
十维交的每个 \(B\in C_{\rm off}\) 提升仍属于 \(E\)。比较第零、
第一行与第 \(u\) 行的块，得到
\(BQ_u^{\mathsf T},BP_u^{\mathsf T}\in C_{\rm off}\)。初等矩阵引理
\[
 BA\text{ 对所有 }B\in C_{\rm off}\text{ 对称}
 \Longrightarrow A\text{ 为标量}
\]
（取 \(B=E_{ij}+E_{ji}\) 比较矩阵元）说明
\(P_u=p_uI,Q_u=q_uI\)。因此
\[
 K=U\otimes C,\qquad
 E\cap\Sym^2K=k(a\odot b)\otimes C_{\rm off},          \tag{SR5}
\]
其中 \(U=\langle a,b\rangle\)、\(a\equiv r_0,b\equiv r_1\pmod t\)。

若某一 Chow 项的本质因子空间至多四维，两项的本质空间之和至多九维；
而 (SR3) 对含 \(M\) 的九平面只允许至多六个矩形，和 (SR5) 的十维
矛盾。故两项的五个因子均独立。投影到 \(a\otimes C\) 后，两项的
共同十维像相同；秩四投影会给同一四维超平面并使二者之和小于十维，
故二者都是图
\[
 L_i=\Gamma(T_i),\qquad T_i:C\to C,\qquad
 \Delta=T_0-T_1\text{ 可逆}.                            \tag{SR6}
\]
以共同 \((a\otimes C)^2\) 分量 \(S\) 标号二次提升；两提升之差
属于 (SR5)，所以
\[
 S\Delta^{\mathsf T}\in C_{\rm off},\qquad
 T_0ST_0^{\mathsf T}=T_1ST_1^{\mathsf T}.              \tag{SR7}
\]
维数相等使第一式的像为整个 \(C_{\rm off}\)。上述矩阵引理给
\(\Delta=\delta I\) 及共同分量空间 \(D=C_{\rm off}\)。置
\(C_0=(T_0+T_1)/2\)，第二式变成
\[
 C_0S+SC_0^{\mathsf T}=0\quad(S\in C_{\rm off}).
\]
再取 \(S=E_{ij}+E_{ji}\) 比较对角元和非对角元，得到 \(C_0=0\)。
于是
\[
 T_0=\frac\delta2I,\qquad T_1=-\frac\delta2I.           \tag{SR8}
\]
若 \(\nu(\delta)>0\)，两个图都专门化到整条第零行；若
\(\nu(\delta)=0\)，两个特殊图仍互不相交；若
\(\nu(\delta)<0\)，除以 \(\delta\) 后二者都专门化到整条第一行。
三者都不可能是 \(M\) 加同一外部行中的一条坐标线。因此 \(d=0\)
也不可能。

只剩 \(d=4\)。此时 \(L_0+L_1\) 是六平面；(SR2)--(SR3) 给
\(\dim(E\cap\Sym^2(L_0+L_1))\le1\)。与 (SR1) 合并即得
\(10-s\le1\)，即 \(s\ge9\)。
\end{proof}

现在排除 \(W_M\)。对 route 4 的四个直和 Chow 块作一个在
\(x_{10},x_{11},x_{20},x_{21}\) 上权重互异、并保持 \(E,M,W_M\)
的一参数行列环面退化。它与原 DVR 参数可在有限多个 Pluecker 权上
用 \(t=s^N,u=s\)、\(N\gg0\) 合并成一条字典序 DVR 弧。四个特殊
五平面因而成为
\[
 M+\langle x_{10}\rangle,\quad M+\langle x_{11}\rangle,\quad
 M+\langle x_{20}\rangle,\quad M+\langle x_{21}\rangle
\]
中的四个，允许重复。每个端点满足
\(E\cap\Sym^2L^*=0\)，故闭商包含与维数给
\(F_i^*=\Sym^2M\)。四个端点只占两个外部行，鸽巢原理给出一对
同行端点；引理~\ref{n5:lem:samerow} 迫使这两个泛纤维二次空间至少
九维相交，与四块直和矛盾。故 \(W_M\) 以及它在
\(t=48,49,50\) 的全部后代均不存在。结合长度二论证，命题
~\ref{n5:prop:orbit1terminal} 的四个临界终端全部排除。

\subsubsection{orbit--13：十四权的统一结构界}

取五格支撑
\[
 L_A=\{00,01,02,10,11\}.
\]
在 \(Q=\Sym^2V/E\) 中，矩形两条对角线给同一个 crossing 权，其余
二次单项的商权互异，故得到十四权宇宙
\[
\begin{aligned}
\mathcal U_A=\{&S_{00},S_{01},S_{02},S_{10},S_{11},
R_{0;01},R_{0;02},R_{0;12},R_{1;01},\\
&C_{01;0},C_{01;1},X_{01;01},X_{01;02},X_{01;12}\}.
\end{aligned}
\]
另一个支撑由行置换得到。环面不变的
\(F\cap\Sym^2L_A\) 在唯一二维 crossing 单项式空间中可以取任意
直线，但商掉 \(E\) 后该权仍只有一维；故 \(\dim q(F)=10\) 时，
\(q(F)\) 正是 \(\mathcal U_A\) 的十元权子集。

这里必须先排除特殊化时的商秩下降。若
\(\dim(E\cap F)=1\)，则 \(q(F)\) 是九维 Chow 商像并包含于特殊
十平面 \(W\)。定理~\ref{n5:thm:one-intersection-flag} 对任意第十
方向直接给出 \(p(W)\le26\)。这已经与本路由要求的
\(p(W)\ge39\) 矛盾。因此
临界分支必有 \(E\cap F=0\) 和 \(q(F)=W\)，上述十四权十子集描述
才被使用；这里没有把像秩在特殊纤维自动保持当作闭条件。

固定三次环面权，把该权的三次单项式作为顶点。每个未选择二次商权
给出的导数条件，或者把一个顶点锚到零，或者给出
\(z_u=\pm z_v\) 的带符号边。选择一个商权等价于删除同标号的边。
因此局部延拓核维数等于不含零顶点且符号相容的连通分量数；所有局部
图至多六顶点。这就是四角四边形、重复行列六边图和 \(K_{3,3}\)
matching 割的统一形式。

写 \(W=C\sqcup X\)，其中 \(C\) 取自前十一条非 crossing 权，
\(X\subset\{X_{01;01},X_{01;02},X_{01;12}\}\)，并令 \(u=|X|\)。
记 \(C\) 中 square、同行边、同列边的数目分别为 \(s,r,c\)。对所选
square \(v\)，记其同行度数、同列度数为 \(a_v,b_v\)，并令
\(\tau\in\{0,1\}\) 表示第零行三条同行边是否全被选入。非 crossing
核由图公式化为
\[
p(C)=5\tau+\sum_v(a_v+1)(b_v+1).
\]
这里 \(a_v\le2,b_v\le1\)，故
\[
(a_v+1)(b_v+1)\le3+a_v+b_v,\qquad
\sum_va_v\le2r,\quad\sum_vb_v\le2c.
\]
若 \(\tau=0\)，立即得到 \(p(C)\le3(s+r+c)\)。若 \(\tau=1\)，则
\(r\ge3\)；当 \(r+c\ge5\) 时同一结论仍由
\(5+3s+2r+2c\le3(s+r+c)\) 得到。当 \(c=0\) 时所有 \(b_v=0\)，故
\(p(C)\le5+3s\le3(s+r)\)。唯一剩余情形是 \((r,c)=(3,1)\)：编号为
\(1\) 的行的同行边未选，唯一同列边只有第零行端点可能比逐点上界 \(3\) 多贡献
\(3\)，于是 \(p(C)\le5+3s+3<3(s+4)\)。因此统一地
\[
p(C)\le3|C|.
\]

再固定一个 crossing \(q=X_{01;ab}\) 和任意
\(D\subset\mathcal U_A\setminus\{q\}\)。加入 \(q\) 对每个变量 \(y\)
只删除权块 \(\operatorname{wt}(q)+\operatorname{wt}(y)\) 中的一条
关系，且二十五个权块互异。若 \(y\notin L_A\) 且该块有重复行或
重复列，则被删边的两个端点各由一个标签在 \(\mathcal U_A\) 外的
square、同行或同列关系锚到零；这些锚不可能被 \(D\) 删除。若三行
三列都互异，则六个顶点是三行到三列的六个双射，九条关系构成
\(K_{3,3}\)；属于 \(\mathcal U_A\) 的边固定行对 \(\{0,1\}\)，因而
构成一个匹配，而 \(K_{3,3}\) 去掉整个匹配后仍是连通六圈。故二十个
\(y\notin L_A\) 的块增量全为零。五个 \(y\in L_A\) 的块各删除一条
关系，核维数各至多增加一。因此
\[
p(D\cup\{q\})-p(D)\le5.
\]
该式对任意 \(D\) 成立，故可以逐次加入 \(X\) 中的 crossing。由
\(|C|=10-u\) 及 \(u\le3\)，
\begin{equation}\label{n5:eq:orbit13}
 |W|=10,\ W\subset\mathcal U_A\quad\Longrightarrow\quad
 p(W)\le3(10-u)+5u\le36<39.
\end{equation}
这是一条不含十权子集分类的结构证明。式~\eqref{n5:eq:orbit13}
排除 orbit--13 及其 \(t=50\) 后代。

\subsubsection{orbit--0：全阶列刚性}

只剩 complete-row 旗标。写 \(A_d,B_d\) 为行、列平方自由 \(d\) 次
空间，并取
\[
 S_0=U_0\otimes B_3,\quad
 U_0=\langle a_{012},a_{013}\rangle,\qquad
 T_0=R_0\otimes B_2,\quad R_0=\partial U_0.
\]
完整旗标关联
\(\mathcal F=\{(S,T):\partial S\subset T\}\) 的线性化方程只有单项锚
和两个系数相等。为使这一点可直接复核，置
\[
\mathcal U=\{012,013\},\qquad
\mathcal R=\{01,02,12,03,13\},
\]
并把两个 Grassmann 图坐标写成
\[
\sigma_{uK}^{vL}\quad
(u\in\mathcal U,\ v\notin\mathcal U,\ |K|=|L|=3),
\qquad
\tau_{rJ}^{sM}\quad
(r\in\mathcal R,\ s\notin\mathcal R,\ |J|=|M|=2).
\]
在 \(\partial_{x_{ia}}(u\otimes b_K)\) 的 \((s,M)\) 商坐标中，
\(\sigma_{uK}^{vL}\) 出现当且仅当
\[
v=s\cup\{i\},\qquad L=M\cup\{a\},
\]
而 \(-\tau_{rJ}^{sM}\) 出现当且仅当
\[
i\in u,\ a\in K,\qquad r=u\setminus\{i\},\
J=K\setminus\{a\}.
\]
故每条方程至多两项且系数为 \(\pm1\)。

若 \(K\ne L\)，取 \(a\in L\setminus K\)。图
\(\mathcal R\) 的三角形恰为 \(\mathcal U\)，故对
\(v\notin\mathcal U\) 可取 \(i\in v\) 使
\(s=v\setminus\{i\}\notin\mathcal R\)。相应方程含
\(\sigma_{uK}^{vL}\) 而不含 \(\tau\)，所以它被锚定。若 \(J\ne M\)，
取 \(a\in M\setminus J\)；由 \(\mathcal R=\partial\mathcal U\)，
可取 \(u=r\cup\{i\}\in\mathcal U\)，再置 \(K=J\cup\{a\}\)。
此时相应方程含 \(\tau_{rJ}^{sM}\)，而因 \(a\in M\) 不可能有
\(L=M\cup\{a\}\)，所以也被锚定。因此所有未锚方向满足
\[
K=L,\qquad J=M.
\]
二部包含图
\[
\binom{[5]}3\longleftrightarrow\binom{[5]}2,\qquad J\subset K,
\]
是连通的，故固定行移动的十份列复制由等系数边粘成一个参数。整个
未锚问题于是降为
\[
2\cdot8+5\cdot5=41
\]
个行侧顶点。逐项使用“目标行三子集删一点落在
\(\binom{[5]}2\setminus R_0\) 即锚定”及其二次对偶判据，唯一自由
分量为
\[
\begin{array}{c|l}
1&012\to014;\ 02\to04,\ 12\to14\\
2&012\to023;\ 12\to23\\
3&012\to024,\ 013\to034;\ 01\to04,\ 12\to24,\ 13\to34\\
4&012\to123;\ 02\to23\\
5&012\to124,\ 013\to134;\ 01\to14,\ 02\to24,\ 03\to34\\
6&013\to014;\ 03\to04,\ 13\to14\\
7&013\to023;\ 13\to23\\
8&013\to123;\ 03\to23.
\end{array}
\]
故切空间维数为八。每个存活变量的列源、列目标相同，而列包含图又把
所有复制等同；所以列环面权为零，列置换只在同一等系数分量内置换。
列群 \(G=T_B\rtimes S_5\) 在切空间上作用平凡。

\begin{lemma}[形式线性化]\label{n5:lem:formal}
若线性约化群 \(G\) 作用在完备 Noether 局部代数
\((R,\mathfrak m)\) 上，且在 \(\mathfrak m/\mathfrak m^2\) 上作用
平凡，则 \(G\) 在 \(R\) 上作用平凡。
\end{lemma}
\begin{proof}
对每个 \(N\)，有限维代数 \(R/\mathfrak m^N\) 有 \(G\)-稳定的
\(\mathfrak m\)-进滤过。其第 \(d\) 个分次商是
\(\Sym^d(\mathfrak m/\mathfrak m^2)\) 的商，因而是平凡表示。
线性约化性使这一有限滤过逐层分裂，所以
\(R/\mathfrak m^N\) 本身是平凡 \(G\)-表示。最后取逆极限
\(R\simeq\varprojlim_N R/\mathfrak m^N\)，得到 \(G\) 在 \(R\) 上作用
恒等。
\end{proof}

应用引理~\ref{n5:lem:formal}，任何专门化到 \((S_0,T_0)\) 的形式弧
都位于列群不动胚。列权分解与 \(S_5\) 传递性给出
\begin{equation}\label{n5:eq:tensorflag}
 S=U_1\otimes B_3,\qquad T=R_1\otimes B_2,\qquad
 \dim U_1=2,\quad\dim R_1=5,\quad\partial S=T.
\end{equation}

\subsubsection{六个 Chow 块的列分离}

在 orbit--0 有 \(h=60\)，故
\(H=F_1\oplus\cdots\oplus F_6\)，每块十维。相应三次空间也直和。
令 \(P_i:T\to F_i,Q_i:S\to J_i\) 为坐标投影，则
\begin{equation}\label{n5:eq:chain}
 \partial_xQ_i(s)=P_i(\partial_xs).
\end{equation}
至少四个饱和块保证每个 \(P_i\) 都满射：给定 \(f_i\)，在另一个饱和
块中消去其共同商像，即得 \(T=E\cap H\) 中以 \(f_i\) 为第 \(i\)
分量的元素。

先固定列三元组 \(I\)。若 \(x\) 的列不在 \(I\)，则
\(\partial_x(U_1\otimes b_I)=0\)，故式~\eqref{n5:eq:chain} 给出
\(\partial_xQ_i(U_1\otimes b_I)=0\)。特征零中这说明
\(Q_i(U_1\otimes b_I)\) 只使用 \(I\) 中三列。记
\(A_a=A\otimes b_a\)，再固定列对
\(\{a,b\}\)。对每个 \(c\notin\{a,b\}\)，由
\(R_1=\partial U_1\)，空间 \(R_1\otimes b_{ab}\) 由
\(U_1\otimes b_{abc}\) 沿第 \(c\) 列的导数组成。因此
\[
P_i(R_1\otimes b_{ab})
\subset\bigcap_{c\notin\{a,b\}}
\Sym^2(A_a\oplus A_b\oplus A_c)
=\Sym^2(A_a\oplus A_b).
\]
还需排除其中的同列平方项；
这是旧稿曾遗漏的步骤。取 \(c\notin\{a,b\}\)、\(u\in U_1\) 及行指标
\(p,r\)，交换混合偏导得到
\begin{equation}\label{n5:eq:mixed}
 \partial_{x_{pa}}P_i((\partial_ru)\otimes b_{ab})
 =\partial_{x_{rc}}P_i((\partial_pu)\otimes b_{bc}).
\end{equation}
右端不含列 \(a\) 一次项，故左端的纯列 \(a\) 平方部分为零；交换
\(a,b\) 排除列 \(b\) 平方。由 \(R_1=\partial U_1\)，
\[
 P_i(R_1\otimes b_{ab})\subset A_aA_b.
\]
不同列对权互异，所以 \(F_i\) 是列环面稳定且无同列项的十维 Chow
二次空间。

若五因子张成维数至多三，则 \(\dim F_i<10\)；若张成四维，则
\(F_i=\Sym^2L_i\) 含平方，也不可能。因此五因子独立。空间
\(F_i=\Span\{\ell_a\ell_b:a<b\}\) 的对偶零化子中，秩一平方恰为
五条坐标线；连通列环面必须逐条固定五个因子线。无同列项又排除两
因子位于同一列，故
\(\ell_{ic}=v_{ic}\otimes b_c\)。最后式~\eqref{n5:eq:chain} 在列三元
组 \(a,b,c\) 上给出非零张量恒等式
\[
 \beta_{ab}\circ\partial=\lambda_{abc}\otimes v_{ic}.
\]
左端与 \(c\) 无关；任意两列之外可选一个不交列对，所以五个
\(v_{ic}\) 全部成比例。于是
\begin{equation}\label{n5:eq:sixblocks}
 F_i=\langle v_i^2\rangle\otimes B_2,\qquad
 J_i=\langle v_i^3\rangle\otimes B_3.
\end{equation}

\subsubsection{任意六个系数的余式至少需十项}

在每个列对块中，模 \(E\) 只保留行对角系数。因此六块共同商像等价于
六个向量 \((v_{i0}^2,\ldots,v_{i4}^2)\) 张成同一条线。逐个缩放
\(v_i\)，再对非零行坐标作一次共同对角缩放，可写成
\(v_i=(\epsilon_{i0}a_0,\ldots,\epsilon_{i4}a_4)\)，其中非零支撑上的
\(\epsilon_{ir}=\pm1\)。完整交空间为
\[
 \widetilde S=
 (A_3\cap\Span\{v_1^3,\ldots,v_6^3\})\otimes B_3.
\]
因为六个立方独立，且
\(\sum c_iv_i^3\in A_3\) 当且仅当 \(\sum c_iv_i=0\)，
\[
 \dim\widetilde S=10(6-\rank\{v_1,\ldots,v_6\}).
\]
左端属于 \(\{20,21,22\}\)，故它等于 20，六个符号向量秩为四，
且行支撑大小为四或五。

令
\[
\Omega=\{(1,z_1,z_2,z_3,z_4):z_i=\pm1\},\qquad
\chi_A(z)=\prod_{i\in A}z_i,
\]
并记 \(RM_{\le d}=\Span\{\chi_A:|A|\le d\}\)。字符正交性给
\[
\dim RM_{\le2}=11,\quad \dim RM_{\le3}=15,\quad
(RM_{\le3})^\perp=\langle\chi_{1234}\rangle.          \tag{**}
\]
若 \(Z\subset\Omega\)，以 \(V_Z\) 表示列为
\((1,z_1,\ldots,z_4)^{\mathsf T}\) 的矩阵，则
\begin{equation}\label{n5:eq:shortening}
\dim\{q\in RM_{\le2}:\operatorname{supp}q\subset Z\}
=|Z|-\rank V_Z.
\end{equation}
事实上，与五个次数至少三的字符正交，在乘以 \(\chi_{1234}\) 后
恰变为 \(V_Zu=0\)。此外，\(r\) 维线性空间中至多有
\(2^{r-1}\) 条射影符号线：选取使坐标投影单射的 \(r\) 个坐标即可。
故六点集 \(\mathcal S\subset\Omega\) 若 \(\rank V_{\mathcal S}=4\)，
则式~\eqref{n5:eq:shortening} 的维数为二，而对每个真子集至多为一。

先处理共同缺一行的情形。若六个 \(v_i\) 都不含 \(x_4\)，则对任意
五次式 \(G(x_0,\ldots,x_3)\)，
\[
R=x_0x_1x_2x_3x_4-G
\]
的六个 \(\partial_i\partial_jR\ (0\le i<j\le3)\) 模掉不含 \(x_4\)
的三次式后给六个不同单项式；另四个
\(\partial_i\partial_4R\) 也是不同单项式。因此行五次式满足
\(\rank C_{2,3}(R)\ge10\)。极化回五个列槽后，十个列三元组给出
彼此支撑互异的十个导数块，总三次导数空间维数至少为一百；而单个
五次 Chow 项的三次导数空间维数至多十。故余式至少需要十个 Chow 项。

以下设满行支撑。对任意标量 \(d_1,\ldots,d_6\)，把六项从行
permanent 中减去：
\[
 R(d)=e_5-\frac1{16}\sum_{z_i\in\mathcal S}
 \chi_{1234}(z_i)d_i\,v(z_i)^5,
\quad v(z)=x_0+\sum_{j=1}^4z_jx_j.
\]
置 \(h(z_i)=1-d_i\)、\(h(z)=1\) 于补集。Glynn Fourier 恒等式与
(**) 给出，在十一维有效二次商上
\[
q\in\ker C_{2,3}(R(d))
\quad\Longleftrightarrow\quad hq\text{ 为常数}.
\]
若该核至少二维，消去常数得到非零 \(q\) 支撑在
\(\mathcal Z=\{z_i:d_i=1\}\)。若某个核元对应非零常数，在
\(\mathcal Z\) 上取值即得矛盾；所以整个核都是支撑在
\(\mathcal Z\) 上的二次 Boolean 函数。式~\eqref{n5:eq:shortening}
迫使 \(\mathcal Z=\mathcal S\)，即
\begin{equation}\label{n5:eq:standardpoint}
d_1=\cdots=d_6=1.
\end{equation}
反之该点的核维正好为二，故行 catalecticant 秩为九。于是非标准
系数时秩至少十，极化到十个列三元组后总三次导数空间至少一百维，
余式不可能由九项表示。

只剩标准点。令 \(\mathcal C=\Omega\setminus\mathcal S\)，则余式为
\[
R_{\mathcal C}=\frac1{16}\sum_{z\in\mathcal C}
\chi_{1234}(z)v(z)^5 .
\]
(**) 还说明任意少于十六个不同 \(v(z)^3\) 线性无关。置
\[
P_3=\Span\{v(z)^3:z\in\mathcal C\},\qquad
R_3=\im C_{2,3}(R_{\mathcal C}).
\]
于是 \(\dim P_3=10,\dim R_3=9\)。由
式~\eqref{n5:eq:shortening}，十个 \(v(z)^2\) 的张成维数为九，故
恰有一条关系
\[
\sum_{z\in\mathcal C}\lambda_zv(z)^2=0,
\]
且其支撑中的符号向量至少张成三维。

若 \(F\in P_3^{(1)}\)，唯一写成
\(F=\sum a_z\otimes v(z)^3\)。前两槽反对称化后，平方关系的一维性
给
\[
a_z\wedge v(z)=\lambda_z\omega.
\]
再与 \(v(z)\) 外积；三个线性无关的支撑向量迫使
\(\omega=0\)。故
\[
P_3^{(1)}=\Span\{v(z)^4:z\in\mathcal C\}.
\]
再令 \(F=\sum c_zv(z)^4\in R_3^{(1)}\)。\(R_3\subset P_3\) 的
超平面方程支撑与 \((\lambda_z)\) 相同；对所有方向收缩得到
\[
\sum_z\eta_zc_zv(z)=0.
\]
其系数矩阵秩至少三，故
\begin{equation}\label{n5:eq:R3prol}
\dim R_3^{(1)}\le10-3=7.
\end{equation}

把标准余式极化回五列。其三次导数空间 \(H_3\) 是十个九维列块的
直和，故 \(\dim H_3=90\)。若一个非零四次列权含重复列，则可选一个
求导列，使导数仍含重复列（例如对型 \(2+1+1\) 取一个单列；其余分拆
更直接），这会产生 \(H_3\) 中不存在的列权。因此延拓只可能有四个
互异列。固定一个
四列块后，沿任一列收缩所得三槽张量都属于对称空间 \(R_3\)；这使原
四槽张量对任意三个槽的置换不变，而这些置换生成 \(\mathfrak S_4\)。
所以该块本身来自一个对称四次式，并恰受
式~\eqref{n5:eq:R3prol} 控制。五个缺列块因而给出
\[
\dim H_3^{(1)}\le5\cdot7=35.
\]
对
\[
 \mathcal K_1:V\otimes H_3\longrightarrow\Lambda^2V\otimes H_2.
\]
其核为 \(H_3^{(1)}\)，从而纯特征零下界为
\[
\rank\mathcal K_1\ge25\cdot90-35=2215.
\]
单个五次 Chow 项在该映射上的秩至多 245：若其三次空间维数至多九，
结论显然；若为十，十个三因子积线性独立。五个四因子积的导数都在
该三次空间中，且把线性关系限制到每个因子超平面便知它们线性独立，
故 prolongation 核至少五维。因而九项总秩至多
\[
9\cdot245=2205<2215.
\]
标准余式也不可能由九项表示。

因此在假设的十五项分解中，取出任意导致 d10 等号的六项后，余下九项
不能表示余式，矛盾。结合表~\ref{n5:tab:routes}，全部 58 态均被排除，
所以
\begin{equation}\label{n5:eq:lower16}
 \ChowRank(\perm_5)\ge16.
\end{equation}

\subsection{Glynn 上界与定理完成}

Glynn 恒等式为
\begin{equation}\label{n5:eq:glynn}
 \perm_5(X)=\frac1{16}
 \sum_{\substack{\epsilon\in\{\pm1\}^5\\ \epsilon_0=1}}
 \left(\prod_{r=0}^4\epsilon_r\right)
 \prod_{c=0}^4\left(\sum_{r=0}^4\epsilon_rx_{rc}\right).
\end{equation}
右端含十六个 Chow 项，故秩至多十六。与
式~\eqref{n5:eq:lower16} 合并，证明定理~\ref{n5:thm:main}。

\subsection{严格依赖边界}

本节的证明先用拆项引理把所有秩至多十五的表达式化为 fixed-six 加
九项余式，因此不以前置 lower--15 为假设。其余依赖为：低秩分类、
一维交旗标定理及其精确模三端点证书、Petersen 乘积影子、双商不等式、
近极大耦合、Borel 闭关联、同时
shifting、层缺陷恒等式给出的十四个 shifted 轮廓、六边父集公式及
四个可见影子结构界、Petersen 纤维轮廓与五平面父集极值、三个终端的局部几何、Boolean
Fourier 刚性及 \(2215>2205\) 的纯 Koszul 估计。orbit--1 的三个
\(W_b\) 由长度二相对局部环排除，第四个 \(W_M\) 由同行全阶赋值引理
排除；四个终端的完备性由命题~\ref{n5:prop:orbit1terminal} 的
五顶点闭式 (T3)、十类 Möbius 项族 (T4) 和八行表 (T11) 纯粹证明。
orbit--13 已由
\(p(C)\le3|C|\) 和单个 crossing 边际至多五的统一结构界处理，不再
含十权子集分类；orbit--0 的一步逆压缩已由 Petersen 偶路引理和五
切片公式处理，不再需要原九十二条森林边。类型 \(1,13\) 的
\(250\) 条 witness 边也已由式~(F1)--(F7) 替换；当前一步逆压缩只含
任意轮廓层缺陷恒等式、三种例外移动模式、五个例外方向和十四个
非统一族状态，再由三个父集结构情形统一排除；旧四十二个方向记录及
十四状态最小值表仅作交叉诊断。

除一维交旗标的独立程序外，参考程序只核对这些结构公式与剩余有限记录。
随机诊断、历史 10GB SAT/DRAT、583-origin 汇总、60,960 旗标闭包、
1096 个标准集、旧有限域 prolongation 真值表、符号六元组轨道、最小
素分解及模素数 Koszul 子式均不在活跃依赖图中。最终内部依赖审计为
\begin{verbatim}
PASS_STRUCTURAL_REDUCTION_LOWER16_INTERNAL_AUDIT
fixed_six_states          = 58
d10_states                = 9
shifted_terminal_profiles = 14
active_1405_profile_enum  = 0
active_fibre_choice_table = 0
active_49_visible_table   = 0
visible_structural_bounds = 4
active_42_direction_table = 0
exceptional_move_schemes  = 3
family_exception_dirs     = 5
nonuniform_family_states  = 14
active_14_state_table     = 0
exceptional_parent_cases  = 3
global_good_planes        = 6
active_witness_edges      = 0
orbit0_witness_edges      = 0
orbit13_structural_bound  = 36
orbit13_subsets_required  = 0
orbit1_terminal_subsets_required = 0
orbit1_terminal_mobius_families  = 10
orbit1_terminal_nonzero_terms    = 50
orbit1_terminal_K4_rows          = 8
orbit1_terminal_max_p     = 36
orbit1_terminal_max_count = 4
orbit1_terminal_QQ_role   = diagnostic
one_intersection_flag_role = theorem_premise
one_intersection_flags     = 886464
one_intersection_max_F3    = 26
old_10GB_assets_required = false
\end{verbatim}
因此本结果不再标作“纯组合”或“program-free”。准确表述是：特征零
代数几何证明，辅以有限组合分类和一个可独立重写、可失败关闭的精确
模三端点证书。附录给出从定义重写该证书及其余诊断检查器所需的全部
坐标、递推、纤维公式与失败条件。
